% Worked solutions for Listas/lista-r9-mergers-acquisitions.md (15 exercises).
% Compile: pdflatex twice; check: pdffonts (Type 1 only, uni = yes).
% Every number is asserted in lista-r9-mergers-acquisitions-solutions_check.py.
\documentclass[11pt,a4paper]{article}
\usepackage[utf8]{inputenc}
\usepackage[T1]{fontenc}
\usepackage{lmodern}   % Type 1 vector fonts (replaces PK bitmaps)
\usepackage{cmap}      % ToUnicode CMap -> copyable text and accents
\usepackage[english]{babel}
\usepackage{amsmath,amssymb,amsfonts}
\usepackage[a4paper,margin=2.2cm]{geometry}
\usepackage{booktabs}
\usepackage{enumitem}
\usepackage{fancyhdr}
\usepackage{xcolor}
\usepackage{microtype}

\pagestyle{fancy}\fancyhf{}
\lhead{\footnotesize Finan\c{c}as II --- Corporativas (MPE Insper)}
\rhead{\footnotesize List R9 (M\&A) --- Worked solutions}
\cfoot{\footnotesize\thepage}
\renewcommand{\headrulewidth}{0.3pt}
\setlength{\headheight}{14pt}
\setlist{topsep=2pt,itemsep=1pt}

\newcommand{\exsec}[1]{\section*{\color{blue!55!black}#1}\addcontentsline{toc}{section}{#1}}
\newcommand{\notation}{\par\noindent\textbf{Notation.}}
\newcommand{\step}[2]{\par\medskip\noindent\textbf{Step #1 --- #2.}~}
\newcommand{\intuition}{\par\medskip\noindent\textbf{Intuition.}~}
\newcommand{\bookans}{\par\medskip\noindent\textbf{Book answer.}~}
\newcommand{\refbib}{\par\smallskip\noindent\textbf{Ref:}~}
\newcommand{\asked}{\par\medskip\noindent\textbf{What is asked.}~}
\newcommand{\approach}{\par\smallskip\noindent\textbf{Approach.}~}
\newcommand{\answer}[1]{\par\smallskip\noindent\fbox{\textbf{Answer #1.}}~}
\newcommand{\pitfalls}{\par\medskip\noindent\textbf{Common mistakes.}}
\newcommand{\E}{\mathbb{E}}

\usepackage[hyphens]{url}
%% STATEMENT-PREAMBLE BEGIN (tools/statements_into_solutions.py)
\usepackage{longtable,array,calc}
\providecommand{\tightlist}{\setlength{\itemsep}{0pt}\setlength{\parskip}{0pt}}
\providecommand{\href}[2]{#2}
\providecommand{\real}[1]{#1}
\makeatletter\@ifundefined{c@none}{\newcounter{none}}{}\makeatother  % pandoc's \LTcaptype{none}
\newcommand{\stmtbegin}{\par\medskip\noindent\textbf{Statement.}\par\nopagebreak\begingroup\small\setlength{\tabcolsep}{3pt}\sloppy}
\newcommand{\stmtend}{\par\endgroup\nopagebreak\noindent\rule{\linewidth}{0.4pt}\par\medskip}
\DeclareUnicodeCharacter{2192}{\ensuremath{\rightarrow}}
\DeclareUnicodeCharacter{2212}{\ensuremath{-}}
\DeclareUnicodeCharacter{2022}{\textbullet{}}
\DeclareUnicodeCharacter{2248}{\ensuremath{\approx}}
\DeclareUnicodeCharacter{2032}{\ensuremath{^{\prime}}}
\DeclareUnicodeCharacter{2081}{\ensuremath{_{1}}}
\DeclareUnicodeCharacter{03B8}{\ensuremath{\theta}}
\DeclareUnicodeCharacter{03BB}{\ensuremath{\lambda}}
\DeclareUnicodeCharacter{03C8}{\ensuremath{\psi}}
\DeclareUnicodeCharacter{03B1}{\ensuremath{\alpha}}
\DeclareUnicodeCharacter{25A0}{\rule{0.7em}{0.7em}}
%% STATEMENT-PREAMBLE END
\begin{document}
\begin{center}{\Large\bfseries Exercise list R9 --- Mergers and acquisitions}\\[4pt]{\large Worked solutions}\end{center}

\noindent Statements are in \url{Listas/lista-r9-mergers-acquisitions.md}; each exercise is cited by number and printed page and its full text is reproduced right before its solution. Short answers: \url{lista-r9-mergers-acquisitions-solutions-short.pdf}. Editions are the local copies: Higgins = Koski, Higgins \& Mitton, \emph{Analysis for Financial Management}, 13th ed. (2023); Grinblatt--Titman, \emph{Financial Markets and Corporate Strategy}, 2nd ed. (2002); Amaro de Matos, \emph{Theoretical Foundations of Corporate Finance} (2001); Tirole, \emph{The Theory of Corporate Finance} (2005). Numbers that illustrate a conceptual answer are labelled ``illustration'' and are not part of the statement.

\tableofcontents

% =============================================================================
\exsec{1. True or false: valuation and acquisitions --- Higgins 9.1, p.~378}
%% BEGIN-STATEMENT 1
\stmtbegin
\textbf{Source:} Higgins (Koski--Higgins--Mitton), Ch. 9, Exercise 9.1, p.~378 (PDF 399) · Answer in
book: yes, pp.~425--426

Is each of the following statements true or false? Explain your answers briefly.

\begin{enumerate}
\def\labelenumi{\alph{enumi}.}
\item
  On average, acquisitions destroy value for shareholders of acquired firms.
\item
  A discounted cash flow valuation of a target company discounts the target's estimated free cash
  flows at the acquirer's cost of capital.
\item
  An acquirer should be willing to pay a higher control premium for a well-managed company than a
  poorly managed one.
\item
  The liquidation value of a company's shares always places a floor under its stock price.
\item
  An unusually low stock price in managements' eyes encourages management to take the company
  private in a management buyout.
\end{enumerate}
\stmtend
%% END-STATEMENT

\asked Five statements to classify as true or false, each with a short reason. (a) Whether
acquisitions destroy value for the shareholders of the \emph{acquired} (target) firm. (b) Which
discount rate a DCF valuation of a target should use. (c) Whether a buyer should pay a larger
control premium for a well-run firm. (d) Whether liquidation value is a floor for the share price.
(e) Whether a share price that management finds too low pushes management toward a buyout.

\approach Higgins' framework for the market for control: the most a buyer can pay is the target's
value under the buyer's control, $V^{\text{new}}$; the stand-alone (minority) value is $V^{\text{sq}}$;
the control premium pays for $V^{\text{new}}-V^{\text{sq}}$, the value of changing how the firm is run.
Discount rates follow the risk of the cash flows (value additivity), and every ``floor'' on a price
requires someone with the power to enforce it.

\notation
\begin{itemize}
  \item $V^{\text{sq}}$: value of the target's equity under current management (the pre-bid market value of a minority stake).
  \item $V^{\text{new}}$: value under the acquirer's management; $V^{\text{new}}-V^{\text{sq}}$: value of control.
  \item $r_T$, $r_A$: costs of capital appropriate to the target's and the acquirer's cash flows.
  \item $L$: liquidation value per share; $p$: market price per share.
\end{itemize}

\step{1}{(a) Target shareholders} We look at who receives the premium. Bids are made well above the
pre-bid price, and the target's price jumps on announcement (Figure~5.3 of Exercise~3: about
$+20$ percentage points on the announcement days, $\approx 28\%$ cumulative). Target shareholders sell
at $V^{\text{sq}}+\text{premium}$.
\answer{(a)} \textbf{False}: targets' shareholders gain on average (returns of 20\% or more around the
announcement); it is the \emph{acquirers'} gains that are doubtful.

\step{2}{(b) Discount rate} We apply value additivity. The value of a stream of cash flows does not
depend on who owns it, so the target's free cash flows $FCF_t^T$ are discounted at the rate that reflects
their own risk:
\begin{align}
  V^T &= \sum_{t\ge 1}\frac{\E[FCF_t^T]}{(1+r_T)^t}. \label{r9e1-dcf}
\end{align}
Using $r_A$ would value a risky target too high when the acquirer is safer ($r_A<r_T$) and too low when it
is riskier.
\answer{(b)} \textbf{False}: discount at the \emph{target's} cost of capital (the risk of the cash flows
being valued), not the acquirer's.

\step{3}{(c) Control premium} We compute the maximum premium as the value of control. Illustration: two
targets each worth $V^{\text{sq}}=100$; under the buyer the poorly run one is worth $140$, the well run one
$105$ (little left to improve). The maximum premium in percent is
\begin{align}
  \frac{V^{\text{new}}-V^{\text{sq}}}{V^{\text{sq}}} &= \frac{(140)-(100)}{100}=40\% \quad\text{(poorly run)}, \nonumber\\
  \frac{V^{\text{new}}-V^{\text{sq}}}{V^{\text{sq}}} &= \frac{(105)-(100)}{100}=5\% \quad\text{(well run)}. \nonumber
\end{align}
\answer{(c)} \textbf{False}: the premium pays for improvements; a well-managed firm offers little to
improve, so the premium should be \emph{smaller}.

\step{4}{(d) Liquidation value} We ask who can force liquidation. A minority shareholder cannot sell
the assets; only the controlling party (or, with dispersed ownership, incumbent management) can. If
they prefer to continue (optimism, private benefits), the price can stay below $L$: illustration, $L=50$
and $p=40$, a gap of $50-40=10$ per share that only a buyer of control can capture.
\answer{(d)} \textbf{False}: liquidation value is a floor only if someone with control is willing to
liquidate; otherwise the shares can trade below it (which is precisely what invites a raider).

\step{5}{(e) MBO motive} Management knows more than the market. Illustration from the book's answer:
price \$30, management's estimate \$80, so buying the firm yields $80-30=50$ per share to the insiders.
\answer{(e)} \textbf{True}: an unusually low price relative to management's private valuation makes a
management buyout attractive, and also creates a conflict of interest (managers gain from a low price).

\intuition All five items read off one identity: price paid $=V^{\text{sq}}+$premium, and the premium
is bounded by the value of control $V^{\text{new}}-V^{\text{sq}}$. Who captures that value depends on
competition among bidders and on information; who can enforce a floor depends on control rights.
\pitfalls
\begin{itemize}
  \item Confusing target and acquirer returns: targets gain, acquirers roughly break even on average.
  \item Using the acquirer's WACC for the target because ``the acquirer is paying''; the discount rate
        belongs to the cash flows.
  \item Treating liquidation value as an arbitrage floor; it is not, because liquidation requires control.
\end{itemize}
\bookans True/false classification matches the book's answer (pp.~425--426): F, F, F, F, T.
\refbib{Higgins Ch.~9, ``The Market for Control'' and ``The Premium for Control'', pp.~354--357; ``The Empirical Evidence'', p.~364; Figure~5.3, p.~169.}

% =============================================================================
\exsec{2. LBO candidates and management buyouts --- Grinblatt \& Titman 20.5 and 20.8, p.~731}
%% BEGIN-STATEMENT 2
\stmtbegin
\textbf{Source:} Grinblatt \& Titman, Ch. 20, Exercises 20.5 and 20.8, p.~731 (PDF 744) · Answer in
book: no

\textbf{20.5.} Leveraged buyouts are observed mainly in industries with relatively stable cash flows
and products that are not highly specialized. Explain why.

\textbf{20.8.} One of the stated benefits of a management buyout is the improvement in management
incentives. In many cases, however, the top managers do not change after the buyout. Explain why.
\stmtend
%% END-STATEMENT

\asked (20.5) Why LBOs cluster in industries with stable cash flows and non-specialized products.
(20.8) Why an MBO improves incentives even though the same managers stay in charge.

\approach An LBO replaces most equity by debt. Debt capacity is limited by expected distress costs,
which depend on cash-flow volatility (risk of default) and on asset specificity (how much value is lost
in liquidation or fire sales). For 20.8 the lever is ownership and the free-cash-flow discipline of
debt (Jensen 1986): the agency problem is about the managers' \emph{stake} and \emph{discretion}, not
their identity.

\notation
\begin{itemize}
  \item $X$: next-period operating cash flow; $D$: face value of debt; riskless debt requires $D\le \min X$.
  \item $s$: managers' fraction of the equity; a \$1 increase in equity value gives managers $s\cdot 1$.
\end{itemize}

\step{1}{(20.5) Stability and debt capacity} We compare how much riskless debt two firms with the same
mean cash flow can carry. Illustration: stable firm $X\in\{90,110\}$, volatile firm $X\in\{40,160\}$,
both with mean $100$. The largest face value repaid in every state is the worst cash flow:
\begin{align}
  D_{\text{stable}} &= 100-10=90, \nonumber\\
  D_{\text{volatile}} &= 100-60=40. \nonumber
\end{align}
The stable firm can borrow more than twice as much without default risk, so the tax shield and the
discipline of debt come with few expected distress costs.

\step{2}{(20.5) Specialization} Specialized assets and products lose value when the firm is distressed:
customers fear for after-sales service and spare parts, suppliers ask for cash, and specialized
assets fetch low prices because few buyers can use them (the Shleifer--Vishny point of Exercise~14).
Expected distress costs are therefore larger for the same default probability.
\answer{(20.5)} Highly levered structures are cheap only when default is unlikely (stable cash flows)
and costly default is mild (generic products and redeployable assets); both reduce the expected costs
of financial distress that limit leverage, while the tax and free-cash-flow benefits remain.

\step{3}{(20.8) Same managers, new incentives} Illustration: before the buyout managers own $s=5\%$;
after it, $s=15\%$ (Kaplan 1989: management ownership roughly triples). For each \$1 of value they create,
\begin{align}
  \text{before: } s\cdot 1 &= 0.05, & \text{after: } s\cdot 1 &= 0.15, & \frac{0.15}{0.05}&=3. \nonumber
\end{align}
In addition, the heavy debt service pre-commits free cash flow, which removes the option to spend it on
negative-NPV projects, and the buyout sponsor sits on a small active board.
\answer{(20.8)} The incentive problem is not that the managers are bad but that they own little and have
discretion over free cash flow. An MBO raises their equity stake and commits cash to debt service and
active owners, so the \emph{same} managers face sharper incentives. Keeping them also keeps their
firm-specific knowledge.

\intuition An LBO trades financial flexibility for discipline. That trade pays where flexibility is
least needed (steady, mature businesses) and where the discipline is most valuable (large free cash flow).
\pitfalls
\begin{itemize}
  \item Saying LBOs need ``high growth''; growth firms need flexibility and have intangible, specific assets.
  \item Thinking an MBO's benefit requires replacing managers; it comes from ownership and debt, not new people.
\end{itemize}
\refbib{Grinblatt \& Titman Ch.~20 (LBOs and MBOs, financial synergies and disciplinary takeovers); Amaro de Matos \S7.2, pp.~187--189.}

% =============================================================================
\exsec{3. Target stock prices around acquisition announcements --- Higgins 5.10, p.~189}
%% BEGIN-STATEMENT 3
\stmtbegin
\textbf{Source:} Higgins (Koski--Higgins--Mitton), Ch. 5, Exercise 5.10, p.~189 (PDF 210) · Answer
in book: no

Suppose in Figure 5.3 that the stock prices of target firms in acquisitions responded to acquisition
announcements over a three-day period rather than almost instantly.

\begin{enumerate}
\def\labelenumi{\alph{enumi}.}
\item
  Would you describe such an acquisition market as efficient? Why or why not?
\item
  Can you think of any trading strategy to take advantage of the delayed price response?
\item
  If you and many others pursued this trading strategy, what would happen to the price response to
  acquisition announcements?
\item
  Some argue that market inefficiencies contain the seeds of their own destruction. In what ways
  does your answer to this problem illustrate the logic of this statement, if at all?
\item
  Immediately after some merger announcements, the stock price of the target firm jumps to a level
  higher than the bid price. Is this proof of market inefficiency? What might explain this price
  pattern?
\end{enumerate}

\paragraph{Context needed to solve it}

\emph{From Higgins Ch. 5, ``How Rapidly Do Stock Prices Adjust to New Information?'', p.~168 (PDF
189):}

Figure 5.3 gives an indication of the speed with which common stock prices adjust to new
information. It is a result of what is known as an event study. In this instance, the researchers
are studying the effect of acquisition offers on the stock prices of target firms. The graph plots
the cumulative percentage change in a target firm's stock price from a period beginning 40 days
before the announcement of the acquisition offer and ending 10 days after. The price changes shown
are averaged over 6,150 acquisition announcements in the U.S. between 1980 and 2008. The date for
each announcement is defined as day 0, so that all announcements can be shown on the same timeline.

An acquisition offer is invariably good news to the target firm's shareholders, because the offer is
at a price well above the prevailing market price of the firm's shares, so we expect to see the
target company's stock price rise after the announcement. The question is: How rapidly? The answer
evident from the graph is: Very rapidly. We see that the stock price drifts upward prior to the
announcement (perhaps due to rumors of a pending announcement), shoots up dramatically on the
announcement day, and then drifts with little direction after the announcement. {[}\ldots{]}

\emph{Figure 5.3, p.~169 (PDF 190), ``Average Stock Price Reaction to Takeover Announcements for
6,150 Target Firms'' (source: Betton, Eckbo, Thompson and Thorburn, Journal of Finance, 2014),
transcribed from the PDF image. Horizontal axis: days relative to takeover announcement (Day 0),
from −45 to 15; vertical axis: cumulative percent change in stock price of takeover target, 0 to 35.
Values read from the curve (approximate):}

{\def\LTcaptype{none} % do not increment counter
\begin{longtable}[]{@{}lllllllllll@{}}
\toprule\noalign{}
Day & −40 & −30 & −20 & −10 & −5 & −1 & 0 & +1 & +5 & +10 \\
\midrule\noalign{}
\endhead
\bottomrule\noalign{}
\endlastfoot
Cumulative change (\%) & 0 & ≈ 0.5 & ≈ 2 & ≈ 4 & ≈ 6 & ≈ 8--9 & ≈ 24 & ≈ 28.5 & ≈ 28 & ≈ 28 \\
\end{longtable}
}

The curve rises slowly from day −40 to day −1, jumps almost vertically between day −1 and day +1,
and is flat (slightly drifting down) from day +1 to day +10.
\stmtend
%% END-STATEMENT

\asked (a) Whether a three-day price response is consistent with efficiency. (b) A trading rule that
would profit from it. (c) What happens to the response if many traders use the rule. (d) How this
shows inefficiencies ``destroying themselves''. (e) Whether a target trading above the bid proves
inefficiency, and what explains it.

\approach Semistrong efficiency: prices reflect all public information at once, so no rule based on
public news earns abnormal returns. A predictable post-announcement drift is a public, exploitable
pattern. The price above the bid is valued as an expectation over possible outcomes (counter-bids),
which is the competition-among-bidders logic of the M\&A models (Exercise~9).

\notation
\begin{itemize}
  \item $CAR_t$: cumulative percentage change in the target's price up to day $t$ (Figure~5.3 values in the list).
  \item Index $100\,(1+CAR_t)$: price relative to day $-40$.
\end{itemize}

\step{1}{(a) Efficiency} From Figure~5.3 the announcement move is from $\approx 8.5\%$ (day $-1$) to
$\approx 28.5\%$ (day $+1$), a jump of $28.5-8.5=20$ points that happens almost instantly. If the same jump were spread
over three days, the day-1 and day-2 returns would be predictable from public news already released.
\answer{(a)} No: a predictable response to public news violates semistrong-form efficiency.

\step{2}{(b) Trading rule} We compute the return of buying at the close of day 0. Suppose one third of the
jump occurs on day 0. The index at the close of day 0 is
\begin{align}
  100+8.5+\frac{20}{3} &= 115.17, \label{r9e3-day0}
\end{align}
and selling at the close of day 2 at $100+28.5=128.5$ yields
\begin{align}
  \frac{128.5}{115.17}-1 &= 11.6\% \label{r9e3-ret}
\end{align}
in two days with no risk specific to the strategy beyond deal risk.
\answer{(b)} Buy target shares as soon as an offer is announced and sell when the price has finished
adjusting (two days later).

\step{3}{(c)--(d) Self-destruction} Purchases at the close of day 0 push the day-0 price up; competition
among traders moves the trade ever earlier until the whole adjustment happens on the announcement, i.e.\
the pattern of Figure~5.3.
\answer{(c)} The response would become (almost) instantaneous; the profit disappears.
\answer{(d)} The profit opportunity attracts the trading that eliminates it; this is why markets with
many informed traders look efficient.

\step{4}{(e) Price above the bid} We value a target share as an expectation over outcomes. Illustration:
bid \$60; with probability $0.3$ a rival bids \$75, with $0.6$ the bid succeeds at \$60, with $0.1$ the deal
fails and the price falls back to \$35:
\begin{align}
  \E[\text{value}] &= (0.3)(75)+(0.6)(60)+(0.1)(35) = 22.5+36+3.5 = 62 > 60. \label{r9e3-exp}
\end{align}
\answer{(e)} No: the market expects a higher (competing or revised) bid; the price above the offer is the
expected value of the share, not a mistake.

\intuition Event studies measure how fast public news enters prices. The jump on day 0 is the
premium being capitalized, and a price above the bid is the market's forecast of a bidding contest.
\pitfalls
\begin{itemize}
  \item Reading the pre-announcement drift (days $-40$ to $-1$) as proof of inefficiency; it is consistent with
        leakage and rumours (and with insider trading, which semistrong efficiency does not rule out).
  \item Concluding in (e) that ``price $>$ bid'' is irrational; it is a probability-weighted value.
\end{itemize}
\refbib{Higgins Ch.~5, ``How Rapidly Do Stock Prices Adjust to New Information?'' and Figure~5.3, pp.~168--169.}

% =============================================================================
\exsec{4. Sources of LBO premiums --- Amaro de Matos 7.2.4 \#3, p.~189}
%% BEGIN-STATEMENT 4
\stmtbegin
\textbf{Source:} Amaro de Matos, Ch. 7, Exercise 7.2.4 \#3, p.~189 (PDF 207) · Answer in book: yes,
p.~271

3. Empirical evidence is consistent in finding large positive premiums for public shareholders and
also finding that these gains are substantially sustained. Explain some of the sources of such large
gains.

\paragraph{Context needed to solve it}

\emph{From Amaro de Matos §7.2.3 (empirical evidence on LBOs), pp.~187--189 (PDF 205--207) --- the
evidence the question summarizes:}

{[}\ldots{]} Kaplan showed that the returns to the prebuyout shareholders are positively and
significantly related to tax benefits from LBOs. The returns to the postbuyout shareholders,
however, are not. They were shown to be positively related to operating income subsequent to the
transaction. This result was seen by Kaplan as evidence that the tax benefits are mostly captured by
the prebuyout shareholders. Besides, the hypothesis that returns to the shareholders are at the
expense of bondholders is rejected. {[}\ldots{]}

The LBOs were described as some sort of concentration of ownership and control. Managers become
owners of the firm, aligning their goals with those of maximizing the firm's value. However, Baker,
Jensen, and Murphy (1988) show that actual executive compensation is very often insufficient to
provide optimal incentives for managers. However, even if that is true, there is still room for
reduction of agency costs in a different way, in particular following Jensen's free cash flow
hypothesis. Increasing debt commits future cash flows, reducing managerial discretion in the
allocation of free cash flow. {[}\ldots{]}

Regarding the control structure of LBOs, Kaplan (1989a) shows that the management postbuyout
ownership increases by about three times the prebuyout values. {[}\ldots{]} More than two-thirds of
the firms disclosed at least one restructuring activity undertaken since the LBO. {[}\ldots{]}
Kaplan (1989a) and Smith (1990) also provide evidence of improved operational results following
LBOs. {[}\ldots{]}

Finally, there is the issue of whether the gains of the prebuyout shareholders come at the expense
of the prebuyout bondholders. Travlos and Cornett (1993) report significant losses at the
statistical level for bondholders and preferred stockholders at the time of the LBO announcement.
These losses are small, however, when compared to the gains to the prebuyout shareholder, thus
failing to provide a satisfactory explanation.
\stmtend
%% END-STATEMENT

\asked Explain where the large and lasting premiums received by pre-buyout public shareholders come from.

\approach A premium is a transfer of part of the value $V^{\text{post}}-V^{\text{pre}}$ to the selling
shareholders. List the sources of that difference: real gains (taxes, incentives, efficiency) and
transfers from other claimants (bondholders, employees) or from information.

\notation
\begin{itemize}
  \item $D$: new permanent debt; $\tau$: corporate tax rate; $\tau D$: present value of a perpetual tax shield (MM with taxes).
  \item $V^{\text{pre}}$, $V^{\text{post}}$: firm value before and after the buyout.
\end{itemize}

\step{1}{Tax benefits} We value the interest tax shield. Illustration: permanent new debt $D=600$ and
$\tau=34\%$:
\begin{align}
  PV(\text{tax shield}) &= \tau D = (0.34)(600)=204. \label{r9e4-ts}
\end{align}
Kaplan finds that pre-buyout shareholders' returns are positively related to these tax benefits.

\step{2}{Agency costs} Management ownership roughly triples (Kaplan 1989a), and debt service commits free
cash flow that managers could otherwise waste (Jensen's free-cash-flow hypothesis); operating
performance improves after the buyout.

\step{3}{Transfers} Bondholders and preferred stockholders lose at the announcement (the new debt is
senior or riskier), and employment grows more slowly; these losses are small relative to the premiums.

\step{4}{Information and efficiency} Managers know more than outside shareholders, so a buyout can
signal undervaluation; and private ownership speeds decisions and avoids disclosure to rivals.
\answer{(\#3)} Tax savings from high leverage; lower agency costs (managerial ownership and the
discipline of debt on free cash flow); wealth transfers from bondholders, preferred stockholders and
employees; managers' private information (undervaluation); efficiency gains of private ownership.

\intuition The premium is sustainable when it rests on real value creation (taxes, incentives,
efficiency); transfers and information effects explain part of it but are too small on their own.
\pitfalls
\begin{itemize}
  \item Listing only taxes: Kaplan shows post-buyout returns follow operating improvements, not taxes.
  \item Treating bondholder losses as the main source; the evidence finds them small.
\end{itemize}
\bookans Matches the book's answer (p.~271): taxes, agency costs, transfers from bondholders, preferred
stockholders and employees, information, efficiency.
\refbib{Amaro de Matos \S7.2.3, pp.~187--189.}

% =============================================================================
\exsec{5. Grossman--Hart with a toehold --- Amaro de Matos 8.1.5 \#1, p.~205}
%% BEGIN-STATEMENT 5
\stmtbegin
\textbf{Source:} Amaro de Matos, Ch. 8, Exercise 8.1.5 \#1, p.~205 (PDF 223) · Answer in book: yes,
p.~271

The next four exercises are related to Grossman and Hart (1990).

1. Show that if the bidder owns some shares of the targeted firm, it is still profitable to bid when
\(p \geq v\).

\paragraph{Context needed to solve it}

\emph{From Amaro de Matos §8.1.1 ``Largely Diffused Ownership'', pp.~193--194 (PDF 211--212):}

A first model of takeovers is the one by Grossman and Hart (1990), who introduce the free-rider
problem. Their idea is to formalize the way in which shareholders free-ride on the bidder potential
improvements to the corporation, reducing the bidder's potential profits to zero and thus making
takeovers unattractive.

Let \(A\) denote the space of possible activities of the management, and \(f(a)\) be the market
value of the shares described as a function of some \(a \in A\). The function \(f\) describes to a
certain extent the ability of the current management under which the maximum possible profit is

\[\max_{a \in A} f(a) = f(a^*).\]

Suppose now that a raider wishes to take over the firm and announces a tender price \(p > f(a^*)\)
at which he/she will buy all available shares. His/her abilities are supposed to be different from
the current management and he/she is supposed to be able to improve the profits by an amount of
\(Z\). In this case, the expected market value is denoted by

\[v = \max_{a \in A} f(a) + Z,\]

where \(Z > 0\). Assume that \(v\) and \(Z\) are known to all the shareholders. Clearly, if
\(v > p > f(a^*)\), the offer looks attractive to the shareholders and to the raider. However, in a
diffusely held firm, such an attempt at control will not succeed. Here is an important assumption to
reach this conclusion:

\textbf{Assumption 8.1.1.} \emph{The ownership of the firm is largely diffused in the sense that
there are a large number of shareholders, each of them holding a very small proportion of the
outstanding shares.}

Therefore, given a tender offer, the impact of each shareholder's tender decision is negligible in
the final result. Under rational expectations, it turns out that the outcome of the bid is not
stochastic. That is to say, the only successful bids are those expected to be successful with a
probability of 1 and all other bids will not happen. This is so, since given a tender offer with
\(p < v\), any shareholder who knows that the offer will succeed prefers to make \(v\) instead of
getting only \(p\). Knowing that his/her decision has no impact on the outcome of the bid, each
shareholder will be better off not tendering (thus receiving \(p\)) than tendering and receiving
\(v\). Therefore, no one tenders, and the bid does not succeed. The alternative would be to tender
at \(p \geq v\), but such conditions are clearly uninteresting to the bidder since his/her profit is
reduced at least to zero.

Takeovers do exist, however. A way to explain them in this context is through different valuations
due to information asymmetries, for instance. In such a case the shareholders would use, say,
\(v_s \neq v\) instead of \(v\), and the bid could be successful at the tender price \(v_s\) at
which the shareholders would be indifferent between tendering and not tendering. If the cost of the
raid is \(c\), the raider's profit in this case would be

\[\Pi_r = v - v_s - c. \qquad (8.1)\]

\emph{(Note: the parenthetical payoffs in ``better off not tendering (thus receiving \(p\)) than
tendering and receiving \(v\)'' are swapped in the book; the argument requires that a non-tendering
shareholder keeps a share worth \(v\) and a tendering one receives \(p\).)}
\stmtend
%% END-STATEMENT

\asked Show that a bidder who already owns shares of the target can profit from a bid at a tender
price $p\ge v$, although without shares such a bid makes a loss.

\approach This is the Grossman--Hart \emph{free-rider problem}: atomistic shareholders tender only if
$p\ge v$, because a non-tendering shareholder keeps a share worth $v$ after a successful bid. With no toehold
the raider then buys the whole improvement at its full value and loses the cost $c$. A toehold earns the
improvement on the shares already owned, which free riding cannot take away (Tirole \S11.5.2.2).

\notation
\begin{itemize}
  \item $q=f(a^\circ)$: value per share under incumbent management (status quo); $v=\max_a f(a)+Z$: value under the raider, $Z>0$; shares normalized to 1.
  \item $p$: tender price per share; $c\ge 0$: cost of the raid; $\theta\in(0,\tfrac12)$: toehold (the book's answer calls it $x$).
\end{itemize}

\step{1}{Why shareholders need $p\ge v$} A shareholder's tender decision does not change the outcome
(Assumption~8.1.1). If the bid succeeds, keeping a share yields $v$ and tendering yields $p$; so tendering
is optimal only if
\begin{align}
  p &\ge v. \label{r9e5-fr}
\end{align}

\step{2}{No toehold} The raider buys all shares at $p$ and owns a firm worth $v$:
\begin{align}
  \Pi_r(\theta=0) &= v-p-c \le v-v-c = -c < 0 \quad\text{when } p\ge v. \label{r9e5-no}
\end{align}

\step{3}{With a toehold} The raider already owns $\theta$ and buys the remaining $1-\theta$ at $p$. After the
bid he owns the whole firm, worth $v$. The book counts his end-of-raid wealth minus the cash spent:
\begin{align}
  \Pi_r &= v-p(1-\theta)-c. \label{r9e5-book}
\end{align}
Setting \eqref{r9e5-book} positive and dividing by $p>0$,
\begin{align}
  v-p+p\theta-c &> 0 \nonumber\\
  p\theta &> p-v+c \nonumber\\
  \theta &> 1+\frac{c-v}{p}. \label{r9e5-cond}
\end{align}
At $p=v$ the threshold is $1+(c-v)/v=c/v$, which is below 1 whenever $c<v$: some toehold suffices.

\step{4}{Incremental profit} Before bidding, the toehold is worth $\theta q$; the raid's \emph{net} gain
subtracts it:
\begin{align}
  \Delta\Pi_r &= v-p(1-\theta)-c-\theta q \nonumber\\
              &= \theta(v-q)-(1-\theta)(p-v)-c. \label{r9e5-inc}
\end{align}
At $p=v$ this is $\theta(v-q)-c$, positive if $\theta>c/(v-q)$. Illustration: $v=100$, $q=80$, $c=1$,
$\theta=0.1$, $p=100$:
\begin{align}
  \Pi_r &= (100)-(100)(1-0.1)-(1) = 100-90-1=9, \nonumber\\
  \Delta\Pi_r &= (0.1)(100-80)-(1) = 2-1 = 1>0, \nonumber
\end{align}
while without a toehold $\Pi_r=100-100-1=-1$. Solving $\Delta\Pi_r=0$ in \eqref{r9e5-inc} for $p$, the highest profitable bid is
$p=v+[\theta(v-q)-c]/(1-\theta)=100+1/0.9=101.11$.
\answer{(\#1)} With a toehold $\theta$, $\Pi_r=v-p(1-\theta)-c>0$ iff $\theta>1+(c-v)/p$ (at $p=v$: $\theta>c/v$);
net of the toehold's initial value the bid pays iff $\theta(v-q)>c+(1-\theta)(p-v)$. The raider profits from the
value created on the shares he already holds.

\intuition Free riding lets the small shareholders capture all the value created on \emph{their} shares;
it cannot touch the value created on the raider's own shares. A toehold therefore restores an incentive
to raid even at a price that leaves the tendering shareholders indifferent.
\pitfalls
\begin{itemize}
  \item Forgetting that tendering shareholders compare $p$ with $v$ (post-raid value), not with $q$.
  \item Counting the raider's profit on the toehold as $\theta v$ instead of $\theta(v-q)$: the toehold was already worth $\theta q$.
\end{itemize}
\bookans Matches the book's answer (p.~271): profit $v-p(1-x)-c>0\iff x>1+(c-v)/p$. The book does not
subtract the toehold's pre-bid value; Step~4 adds that refinement.
\refbib{Amaro de Matos \S8.1.1, pp.~193--194; Tirole \S11.5.2.2, p.~433.}

% =============================================================================
\exsec{6. News Corp.\ and Dow Jones --- Higgins 9.3, p.~379}
%% BEGIN-STATEMENT 6
\stmtbegin
\textbf{Source:} Higgins (Koski--Higgins--Mitton), Ch. 9, Exercise 9.3, p.~379 (PDF 400) · Answer in
book: yes, p.~426

In July 2007, News Corp.~entered into an agreement to purchase all of the outstanding shares of Dow
Jones and Company for \$60 per share. Immediately prior to the News Corp.~bid, the shares of Dow
Jones traded at \$33 per share. The number of outstanding shares at the time of the announcement was
82 million. The book value of interest-bearing liabilities on the balance sheet of Dow Jones was
\$1.46 billion.

\begin{enumerate}
\def\labelenumi{\alph{enumi}.}
\item
  Estimate the cost of this acquisition to the shareholders of News Corp.
\item
  What value did News Corp.~place on the control of Dow Jones and Company?
\end{enumerate}
\stmtend
%% END-STATEMENT

\asked (a) The cost of the acquisition to News Corp.'s shareholders. (b) The value News Corp.\ placed
on control of Dow Jones.

\approach Buying all the shares of a company means buying its assets and \emph{assuming} its
liabilities, so the cost of the acquisition is the price paid for equity plus the (market value,
approximated by book value) of the interest-bearing debt taken over. The value of control is the premium
over the pre-bid (minority) market value.

\notation
\begin{itemize}
  \item $p_b=\$60$: bid per share; $p_0=\$33$: pre-bid price; $n=82$ million shares.
  \item $D=\$1.46$ billion: interest-bearing liabilities (book value, used as proxy for market value).
\end{itemize}

\step{1}{(a) Equity paid} We multiply the bid by the shares outstanding:
\begin{align}
  p_b n &= (60)(82\text{ million}) = \$4.92\text{ billion}. \label{r9e6-eq}
\end{align}

\step{2}{(a) Add the assumed debt} The buyer becomes responsible for Dow Jones's debt:
\begin{align}
  \text{Cost} &= p_b n + D = (4.92)+(1.46) = \$6.38\text{ billion}. \label{r9e6-cost}
\end{align}
\answer{(a)} About \$6.38 billion (an estimate because book debt approximates market debt).

\step{3}{(b) Value of control} We compute the premium over the pre-bid price, per share and in total:
\begin{align}
  (p_b-p_0)\,n &= (60-33)(82\text{ million}) = (27)(82\text{ million}) = \$2.214\text{ billion}, \label{r9e6-ctrl}\\
  \frac{p_b-p_0}{p_0} &= \frac{27}{33} = 81.8\%. \nonumber
\end{align}
\answer{(b)} About \$2.21 billion, an 81.8\% premium over the pre-bid equity value of $33\times 82=\$2.706$ billion.

\intuition The pre-bid price is the value of a minority stake under current management; News Corp.\ paid
\$2.21 billion more because it expected to run Dow Jones differently (or to enjoy private benefits).
\pitfalls
\begin{itemize}
  \item Reporting only the \$4.92 billion paid for equity: the assumed debt is part of the price.
  \item Computing the control value from total cost (\$6.38 billion) instead of the premium over the pre-bid equity value.
\end{itemize}
\bookans Matches the book's answer (p.~426): \$6.38 billion and \$2.21 billion.
\refbib{Higgins Ch.~9, ``The Premium for Control'', pp.~354--357.}

% =============================================================================
\exsec{7. Procureps: EPS dilution as a merger criterion --- Higgins 9.9, p.~381}
%% BEGIN-STATEMENT 7
\stmtbegin
\textbf{Source:} Higgins (Koski--Higgins--Mitton), Ch. 9, Exercise 9.9, p.~381 (PDF 402) · Answer in
book: yes, p.~427

Procureps, Inc.~(P) is considering two possible acquisitions, neither of which promises any
enhancements or synergistic benefits. V1 is a poorly performing firm in a declining industry with a
price-to-earnings ratio of 8 times. V2 is a high-growth technology company with a price-to-earnings
ratio of 35 times. Procureps is interested in making any acquisition that increases its current
earnings per share. All of Procureps's acquisitions are exchange-of-share mergers.

\begin{enumerate}
\def\labelenumi{\alph{enumi}.}
\item
  Calculate the maximum percentage premium Procureps can afford to pay for V1 and V2 by replacing
  the question marks in the following table.
\item
  What do your answers to part (a) suggest about the wisdom of using ``avoid dilution in earnings
  per share'' as a criterion in merger analysis?
\end{enumerate}

{\def\LTcaptype{none} % do not increment counter
\begin{longtable}[]{@{}llllll@{}}
\toprule\noalign{}
Company & P & V1 & P + V1 & V2 & P + V2 \\
\midrule\noalign{}
\endhead
\bottomrule\noalign{}
\endlastfoot
Earnings after tax (\$ millions) & 2 & 1 & 3 & 1 & 3 \\
Price-to-earnings ratio (X) & 30 & 8 & & 35 & \\
Market value of equity (\$ millions) & ? & ? & & ? & \\
Number of equity shares (millions) & 1 & 1 & ? & 1 & ? \\
Earnings per share (\$) & 2 & 1 & 2 & 1 & 2 \\
Price per share (\$) & ? & ? & & ? & \\
Maximum new shares issued (millions) & & & ? & & ? \\
Value of new shares issued (\$ millions) & & & ? & & ? \\
Maximum acquisition premium (\%) & & & ? & & ? \\
\end{longtable}
}
\stmtend
%% END-STATEMENT

\asked (a) Fill in the table: market values, prices, the maximum number of shares Procureps (P) can
issue without diluting its EPS of \$2, their value and the implied maximum premium for V1 and for V2.
(b) What this says about the ``avoid EPS dilution'' rule.

\approach In an exchange-of-shares merger with no synergies, the combined earnings are the sum of the
earnings. Keeping EPS at \$2 caps the combined share count; the extra shares are worth P's price; the
premium compares that value with the target's market value. The rule depends only on P/E ratios, not on
value creation.

\notation
\begin{itemize}
  \item $E_P=2$, $E_{V1}=E_{V2}=1$: earnings (\$ million); P/E: $30$, $8$, $35$; one million shares each.
  \item $n^{\max}$: maximum combined shares with EPS $\ge 2$; $\Delta n=n^{\max}-1$: new shares; $p_P$: P's share price.
\end{itemize}

\step{1}{Market values and prices} Market value $=$ P/E $\times$ earnings, and price $=$ value $/$ shares:
\begin{align}
  MV_P &= (30)(2)=60, & p_P&=60/1=\$60, \nonumber\\
  MV_{V1} &= (8)(1)=8, & p_{V1}&=\$8, \nonumber\\
  MV_{V2} &= (35)(1)=35, & p_{V2}&=\$35. \nonumber
\end{align}

\step{2}{Maximum shares} Combined earnings are $2+1=3$ in either merger. EPS stays at 2 if
\begin{align}
  \frac{3}{n^{\max}} &= 2 \;\Rightarrow\; n^{\max}=\frac{3}{2}=1.5 \text{ million}, \qquad \Delta n = 1.5-1 = 0.5\text{ million}. \label{r9e7-n}
\end{align}

\step{3}{Value of the new shares and premium} At P's price the 0.5 million new shares are worth
$(0.5)(60)=\$30$ million. The maximum premium over each target's market value is
\begin{align}
  \text{V1: }\frac{30}{8}-1 &= 3.75-1 = 275\%, \label{r9e7-v1}\\
  \text{V2: }\frac{30}{35}-1 &= 0.857-1 = -14.3\%. \label{r9e7-v2}
\end{align}

\begin{center}
\begin{tabular}{lccccc}
\toprule
Company & P & V1 & P+V1 & V2 & P+V2\\
\midrule
Earnings (\$m) & 2 & 1 & 3 & 1 & 3\\
P/E & 30 & 8 & & 35 & \\
Market value (\$m) & 60 & 8 & & 35 & \\
Shares (m) & 1 & 1 & 1.5 & 1 & 1.5\\
EPS (\$) & 2 & 1 & 2 & 1 & 2\\
Price (\$) & 60 & 8 & & 35 & \\
Max.\ new shares (m) & & & 0.5 & & 0.5\\
Value of new shares (\$m) & & & 30 & & 30\\
Max.\ premium (\%) & & & 275 & & $-14$\\
\bottomrule
\end{tabular}
\end{center}
\answer{(a)} Procureps can pay up to a 275\% premium for V1 but must buy V2 at a 14\% \emph{discount} to
avoid EPS dilution.

\step{4}{(b) What the rule does to value} With no synergies, value is conserved. Buying V2 at its market
price issues $35/60=0.583$ million shares; EPS falls to $3/1.583=\$1.89$, but the share price is unchanged:
\begin{align}
  \frac{60+35}{1+0.583} &= \frac{95}{1.583} = \$60. \nonumber
\end{align}
Buying V1 at the ``allowed'' 275\% premium keeps EPS at \$2 but transfers $30-8=\$22$ million to V1's
holders: P's price falls to $(60+8)/1.5=\$45.33$, a loss of \$14.67 per share.
\answer{(b)} The criterion is shortsighted: it favours low-P/E, low-growth targets and penalizes
high-growth ones, and it can approve value-destroying deals while rejecting value-neutral ones. Value, not
EPS, is the test.

\intuition EPS accretion is an artefact of exchanging high-P/E paper for low-P/E earnings; the market
prices the combined firm at a blended P/E, so EPS moves while value does not.
\pitfalls
\begin{itemize}
  \item Valuing the new shares at the target's price instead of the acquirer's.
  \item Reading EPS dilution as value destruction (the V2 case shows EPS down, price unchanged).
\end{itemize}
\bookans Matches the book's answer (p.~427): 275\% and $-14\%$.
\refbib{Higgins Ch.~9, ``Avoiding Dilution in Earnings per Share'', p.~362.}

% =============================================================================
\exsec{8. Hansen: cash versus stock with a uniform target value --- Amaro de Matos 8.2.6 \#3, p.~223}
%% BEGIN-STATEMENT 8
\stmtbegin
\textbf{Source:} Amaro de Matos, Ch. 8, Exercise 8.2.6 \#3, p.~223 (PDF 241) · Answer in book: yes,
p.~276

3. In the model of Hansen (1987), let \emph{v} be uniformly distributed over (0, 100) and take
\emph{w}(\emph{v}) = 1.5\emph{v}. Show that no mutually beneficial cash trade exists, but that a
stock deal will work if \(x < 150\).

\paragraph{Context needed to solve it}

\emph{From Amaro de Matos §8.2.2 ``Choosing the Means of Payment'', pp.~207--209 (PDF 225--227):}

In a model developed by Hansen (1987), payment is made either in cash or in stock of the merged
corporation. The basis of the mechanism that makes management and bidder prefer one or the other
(either cash or stock) is asymmetry of information.

The acquirer owns a firm of value \(x\), which is known to both the acquirer and the target's
management. The target firm is worth \(v\) to itself, but has a value \(w(v) \geq v\) to the
acquirer. The value of a merger will therefore be \(x + v\). The asymmetry enters assuming that the
bidder does not know the value of the firm \(v\), but the current management does. Then, the target
will never accept an offer below that value in a cash trade. This creates an adverse selection
problem, against which the acquirer can protect himself/herself, defining his/her optimal strategy
not on unconditional expected profit, but on expected profit conditional on the offer being
accepted.

The acquirer takes \(v\) as a random variable with support on the closed interval \([v_l, v_h]\).
For a cash offer \(C\), the expected profit of the raider is an expectation conditional on the
acceptance of the offer---that is to say, conditional on \(v_l < v < C\). If \(f(v)\) denotes the
density of probability of \(v\),

\[E(\Pi_r \mid C) = \int_{v_l}^{C} [w(v) - C] f(v)\, dv.\]

Denoting by \(F(v)\) the distribution of \(v\) and by \(E(w \mid C)\) the expected value of the
target to the acquirer conditional on sale,

\[E(w \mid C) = \frac{1}{F(C)} \int_{v_l}^{C} w(v) f(v)\, dv,\]

the expected conditional wealth gain can be rewritten as

\[E(\Pi_r \mid C) = F(C)[E(w \mid C) - C].\]

The optimal cash offer \(C\) must satisfy the first-order condition

\[w(C) - C - \frac{F(C)}{f(C)} = 0\]

together with the obvious restrictions of positive expected profit

\[E(\Pi_r \mid C) \geq 0\]

and \(v_l \leq C \leq v_h\). If \(E(\Pi_r \mid C) < 0\), which is to say when \(C > E(w \mid C)\),
or if the solution is at \(v_l\), no cash trade will take place.

But if the bidder offers stock of the merged firm (say, a proportion \(\beta\)), trade may still
occur. From the point of view of the acquirer, such a trade is in the interest of the target with
certainty if the final value held by the owners of the target is greater than or equal to \(v\) for
whatever current value \(v \in [v_l, v_h]\). It will be in the interest of the target with
probability \(F(v^*)\) if the same inequality holds for any \(v \in [v_l, v^*]\). If the management
of the target knows the function \(w\), this means that with probability \(F(v^*)\) trade is good
for the target if

\[\beta[x + w(v)] \geq v \quad \forall v \in [v_l, v^*]\]

or

\[\beta \geq \frac{v}{x + w(v)} \quad \forall v \in [v_l, v^*], \qquad (8.11)\]

which is equivalent to

\[\beta \geq \max_v \frac{v}{x + w(v)}.\]

Making Hansen's additional assumption that

\[\frac{dw}{dv} < \frac{x + w(v)}{v}\]

the preceding condition becomes simply

\[\beta[x + w(v^*)] \geq v^*.\]

More easily, the acquirer will be interested only if the value held by him/her after the transaction
is greater than the current value. In other words,

\[(1 - \beta)[x + E(w \mid v^*)] \geq x. \qquad (8.12)\]

These two conditions hold simultaneously if

\[\frac{E(w \mid v^*)}{\beta} - E(w \mid v^*) \geq x \geq \frac{v^*}{\beta} - w(v^*),\]

implying that

\[x \leq \frac{w(v^*) - v^*}{v^* - E(w \mid v^*)}\, E(w \mid v^*).\]
\stmtend
%% END-STATEMENT

\asked Show that with $v\sim U(0,100)$ and $w(v)=1.5v$ no cash offer benefits both sides, while an
offer of a fraction $\beta$ of the merged firm's shares is accepted by every target type and still
benefits the acquirer whenever its own value $x$ is below 150.

\approach Hansen's model is a lemons problem: the target knows $v$, the acquirer does not. A cash offer
$C$ is accepted only by targets with $v\le C$ (adverse selection), so the acquirer's gain is computed
\emph{conditional on acceptance}. A stock offer makes the target's payment contingent on what the
target is worth (the target receives a share of $x+w(v)$), which weakens the adverse selection: the
means of payment matters under asymmetric information.

\notation
\begin{itemize}
  \item $v\in(0,100)$: target's stand-alone value, privately known to the target; $F(v)=v/100$, $f(v)=1/100$.
  \item $w(v)=1.5v$: target's value to the acquirer; $x>0$: acquirer's value, publicly known.
  \item $C$: cash offer; $\beta\in(0,1)$: fraction of the merged firm offered; $v^*$: highest accepting type.
  \item $E(w\mid C)=\E[w(v)\mid v\le C]$: value to the acquirer of the targets that accept.
\end{itemize}

\step{1}{Who accepts a cash offer} The target sells for $C$ only if $C\ge v$, so the accepted types are
$v\in(0,C]$ with probability $F(C)=C/100$ (for $C\le 100$).

\step{2}{Value of what the acquirer gets} Conditional on $v\le C$, $v$ is uniform on $(0,C)$ with mean
$C/2$, so
\begin{align}
  E(w\mid C) &= 1.5\,\E[v\mid v\le C] = 1.5\cdot\frac{C}{2} = 0.75\,C. \label{r9e8-Ew}
\end{align}

\step{3}{Cash gain} Substituting \eqref{r9e8-Ew} into $E(\Pi_r\mid C)=F(C)[E(w\mid C)-C]$,
\begin{align}
  E(\Pi_r\mid C) &= \frac{C}{100}\,\big[(0.75C)-C\big] \nonumber\\
                 &= -\frac{0.25\,C^2}{100} < 0 \quad\text{for every } C\in(0,100]. \label{r9e8-cash}
\end{align}
(At $C=100$, where every type accepts, the loss is $-25$.) The acquirer always pays $C$ for assets worth
$0.75C$ to it: although $w(v)=1.5v>v$ creates gains from trade for every type, no cash offer is
profitable.
\answer{(cash)} No mutually beneficial cash trade exists: $E(\Pi_r\mid C)=-0.0025\,C^2<0$ for all $C$.

\step{4}{Who accepts a stock offer} A target of type $v$ holding $\beta$ of the merged firm gets
$\beta[x+w(v)]$; it accepts iff
\begin{align}
  \beta(x+1.5v) &\ge v \iff \beta \ge \frac{v}{x+1.5v}. \label{r9e8-acc}
\end{align}
The right-hand side increases in $v$: its derivative is
\begin{align}
  \frac{d}{dv}\frac{v}{x+1.5v} &= \frac{(x+1.5v)-1.5v}{(x+1.5v)^2} = \frac{x}{(x+1.5v)^2} > 0, \nonumber
\end{align}
which is Hansen's condition $dw/dv=1.5<(x+1.5v)/v$. So the hardest type to please is $v=100$, and every
type accepts iff
\begin{align}
  \beta(x+150) &\ge 100 \iff \beta \ge \frac{100}{x+150}. \label{r9e8-bmin}
\end{align}

\step{5}{Acquirer's participation} If every type accepts, the acquirer gets $w(v)$ with mean
$\E[1.5v]=1.5\times 50=75$, and keeps $1-\beta$ of the merged firm. Condition \eqref{r9e8-acc} of the
model, (8.12), is
\begin{align}
  (1-\beta)(x+75) &\ge x \iff \beta \le 1-\frac{x}{x+75} = \frac{75}{x+75}. \label{r9e8-bmax}
\end{align}

\step{6}{Both conditions} A $\beta$ exists iff the lower bound \eqref{r9e8-bmin} does not exceed the
upper bound \eqref{r9e8-bmax}. Cross-multiplying (both denominators positive):
\begin{align}
  \frac{100}{x+150} &\le \frac{75}{x+75} \nonumber\\
  100(x+75) &\le 75(x+150) \nonumber\\
  100x+7500 &\le 75x+11250 \nonumber\\
  25x &\le 3750 \nonumber\\
  x &\le 150. \label{r9e8-x}
\end{align}
With $x<150$ the inequality is strict, so some $\beta$ gives the acquirer a strict gain. Example $x=100$:
$\beta\in[100/250,\,75/175]=[0.4,\,0.4286]$; at $\beta=0.4$ the acquirer's gain is
$(1-0.4)(100+75)-100=105-100=5$.

\step{7}{Necessity (partial-acceptance stock offers)} An offer accepted only by $v\le v^*$ sets
$\beta=v^*/(x+1.5v^*)$ and yields conditional value $E(w\mid v^*)=0.75v^*$, so the gain is
$F(v^*)\{(1-\beta)(x+0.75v^*)-x\}$. Since $1-\beta=(x+0.5v^*)/(x+1.5v^*)$,
\begin{align}
  (1-\beta)(x+0.75v^*)-x &= \frac{(x+0.5v^*)(x+0.75v^*)-x(x+1.5v^*)}{x+1.5v^*} \nonumber\\
  &= \frac{x^2+1.25xv^*+0.375v^{*2}-x^2-1.5xv^*}{x+1.5v^*} \nonumber\\
  &= \frac{v^*(0.375v^*-0.25x)}{x+1.5v^*}, \nonumber
\end{align}
positive iff $v^*>\tfrac{2}{3}x$. As $v^*\le 100$, some stock offer gains iff $x<150$: the bound is also
necessary.
\answer{(stock)} Any $\beta\in[100/(x+150),\,75/(x+75)]$ is accepted by all types and benefits the acquirer;
the interval is non-empty iff $x\le 150$ (strict gain for $x<150$), and no stock offer gains for $x\ge 150$.

\intuition Cash is a fixed payment: only targets worth less than the offer accept, so the acquirer buys
lemons. Stock pays each target in proportion to the merged value, which rises with $w(v)$; part of the
payment adjusts with the target's quality. The adjustment is diluted when the acquirer is large, because
then $\beta(x+w)$ moves little with $v$; hence the upper bound on $x$.
\pitfalls
\begin{itemize}
  \item Evaluating the cash offer at the unconditional mean $\E[w]=75$ instead of the mean among
        targets that accept ($0.75C$): the winner's curse.
  \item Checking the target's condition at the mean $v$ instead of at the highest accepting type.
  \item Reading the book's ``$p\le 75/(x+75)$'' literally; the $p$ is a misprint for $\beta$.
\end{itemize}
\bookans Matches the book's answer (p.~276): the acquirer gets $0.75C$ for $C$; stock works when
$\beta(x+150)\ge 100$ and $(1-\beta)(x+75)\ge x$, consistent iff $x\le 150$.
\refbib{Amaro de Matos \S8.2.2, pp.~207--210; Hansen (1987).}

% =============================================================================
\exsec{9. Fishman (1988): why the preemptive signal is not fully revealing --- Amaro de Matos 8.2.6 \#1, p.~223}
%% BEGIN-STATEMENT 9
\stmtbegin
\textbf{Source:} Amaro de Matos, Ch. 8, Exercise 8.2.6 \#1, p.~223 (PDF 241) · Answer in book: yes,
p.~276

1. Explain why the signaling equilibrium suggested by Fishman (1988) is not fully revealing. If
\(v_1\) could be privately observed by bidder 2, how would that change the signaling equilibrium?

\paragraph{Context needed to solve it}

\emph{From Amaro de Matos §8.2.1 ``Competition between Bidders'', pp.~206--207 (PDF 224--225):}

The first model to be described here along these lines is the one by Fishman (1988) in which there
are two bidders (say, bidder 1 and bidder 2), who make successive bids until only one remains. The
remaining bid is accepted by the target, provided that its value is greater than or equal to the
value \(v\), which is assumed to be reasonable by the target. Otherwise, there is no takeover. Here,
\(v\) is taken as the current market value of the target, assuming that there is no takeover.

Let bidder 1 be the first to make an offer. The costs \(c_i\) for \(i = 1, 2\) of the bid are the
costs to the bidders of private observation of their own valuations \(v_i\), \(i = 1, 2\).
Furthermore, it is assumed that the valuations are random variables such that \(E(\tilde v_i) < v\).
This means that bidders never make a bid without first observing whether the realization of their
valuation is above \(v\). If \(p\) is the initial offer made by bidder 1, it must be such that
\(v \leq p \leq v_1\) and at the same time makes competition unattractive to bidder 2. In other
words, the expected payoff from competing to bidder 2 should be nonpositive.

The payoffs can be expressed simply in terms of a variable \(d\) that takes the value 0 if bidder 2
does not compete and takes 1 otherwise. For any two random realizations of \(v_1\) and \(v_2\) the
payoff for bidder 1 is

\[\Pi_1(v_1, v_2, p, d) = v_1 - (1 - d)p - d \min[\max(p, v_2), v_1] - c_1 \qquad (8.9)\]

and for bidder 2 is

\[\Pi_2(v_1, v_2, p, d) = d[(v_2 - \min(v_2, v_1) - c_2]. \qquad (8.10)\]

Bidder 2 does not have access to \(v_1\) as bidder 1 does not know \(v_2\). But both of them know
the densities of probabilities \(f_i(v_i)\), which are assumed to be strictly positive in a closed
interval \([l, h]\) and zero elsewhere. It is clear that if \(p\) is very large, it will discourage
bidder 2 from making a higher bid. If it is not too costly, the ideal for bidder 1 is to find the
threshold above which bidder 2 will not even incur the cost of being informed of his/her valuation
\(v_2\). Denoting this threshold by \(r\), it must be such that the expected payoff of bidder 2
conditional on \(v_1 \geq r\) is zero:

\[\frac{\int_r^h E_{v_2} \Pi_2(v_1, \tilde v_2, p, 1) f_1(v_1)\, dv_1}{\int_r^h f_1(v_1)\, dv_1} = 0.\]

For any bid \(v'\) higher than \(r\), the expected payoff of bidder 2 will be negative. On the other
hand, bidder 1 will be able to preempt only if \(v_1 > r\). In the event that \(v_1 \leq r\),
his/her best alternative is to offer \(v\). Therefore, if he/she preempts, the gain he/she gets for
doing so at an offer price \(p'\) is

\[E_{v_2}\Pi_1(v_1, \tilde v_2, p', 0) - E_{v_2}\Pi_1(v_1, \tilde v_2, v, 1) = E_{v_2} \min[\max(v, \tilde v_2), v_1] - p'.\]

Let the highest price that bidder 1 is willing to pay be denoted by \(\bar p(v_1)\). It has to be
such that the preceding gain is still positive. Noticing that the gain is decreasing in \(p'\), it
follows that

\[\bar p(v_1) = E_{v_2} \min[\max(v, \tilde v_2), v_1].\]

It can be shown that it is an equilibrium to bidder 1 to offer \(\bar p(r)\) if \(v_1 \geq r\) and
\(v\) otherwise, whereas for bidder 2 it is an equilibrium not to bid in the first case and to bid
in the second. It is a signaling equilibrium, although it is not fully revealing.

This model implies several things. It is clear that the expected payoff of bidder 1 increases with
the cost \(c_2\), since it is easier to preempt the entry of bidder 2 (\(r\) decreases). If the
first bid can be lower, the expected payoff of the target's shareholders will also be lower.
{[}\ldots{]}
\stmtend
%% END-STATEMENT

\asked Why bidder 1's opening bid in Fishman's equilibrium reveals only a range of $v_1$, and how the
equilibrium would change if bidder 2 could observe $v_1$.

\approach Preemptive bidding is signalling: a high opening bid deters bidder 2 from paying $c_2$ to
enter an auction. Deterrence only requires bidder 2 to believe that $v_1$ is high enough on average
(above the threshold $r$); once that belief is induced, a higher type gains nothing by bidding more.
Compare with the complete-information benchmark where no signal is needed.

\notation
\begin{itemize}
  \item $v$: target's stand-alone value; $v_i$: bidder $i$'s valuation, density $f_i$ on $[l,h]$; $c_2$: bidder 2's cost of learning $v_2$.
  \item $g(v_1)\equiv\E_{v_2}[\Pi_2(v_1,\tilde v_2,\cdot,1)]+c_2=\E[\max(\tilde v_2-v_1,0)]$: bidder 2's gross expected gain from an auction against $v_1$, decreasing in $v_1$.
  \item $r$: preemption threshold, $\E[g(\tilde v_1)\mid \tilde v_1\ge r]=c_2$; $w$: complete-information entry cutoff, $g(w)=c_2$.
  \item $\bar p(v_1)=\E_{v_2}\min[\max(v,\tilde v_2),v_1]$: the highest preemptive bid type $v_1$ is willing to make.
\end{itemize}

\step{1}{Bidder 2's payoff} From (8.10) with $d=1$, $\Pi_2=v_2-\min(v_2,v_1)-c_2=\max(v_2-v_1,0)-c_2$:
in the ascending auction bidder 2 wins when $v_2>v_1$ and pays $v_1$. Hence its expected payoff is
$g(v_1)-c_2$, and $g$ falls as $v_1$ rises.

\step{2}{What the equilibrium reveals} Bidder 1 offers $\bar p(r)$ if $v_1\ge r$ and $v$ if $v_1<r$.
Bidder 2 sees one of two messages: ``$v_1\in[v,r)$'' or ``$v_1\ge r$''. All types above $r$ make the
\emph{same} bid. They do not separate because bidder 2's decision depends only on its belief about
$v_1$, and the pooled belief $\E[g\mid v_1\ge r]=c_2$ already makes it stay out; bidding more costs a type
above $r$ money and changes nothing.
\answer{(why)} The equilibrium is a two-message (partition) equilibrium: bidder 2 learns only whether
$v_1<r$ or $v_1\ge r$, not $v_1$ itself, so it is not fully revealing.

\step{3}{Complete information} If bidder 2 observed $v_1$, it would enter iff $g(v_1)\ge c_2$, i.e.\
iff $v_1<w$ with $g(w)=c_2$. Bidder 1 would no longer need a costly signal; he would open at $v$ (the
lowest acceptable price) and let bidder 2 decide on the true $v_1$. Because $g$ is decreasing and
$\E[g\mid v_1\ge r]=c_2$, the value at the bottom of the pool is above the pool average:
$g(r)>c_2=g(w)$, hence $w>r$.

\step{4}{Illustration} Take $v_1,v_2\sim U[0,1]$, $v=0.6$ (so $\E[\tilde v_i]=0.5<v$) and $c_2=0.02$.
Then $g(v_1)=\int_{v_1}^1(t-v_1)\,dt=(1-v_1)^2/2$.
\begin{align}
  \text{cutoff } w:&\quad \frac{(1-w)^2}{2}=0.02 \;\Rightarrow\; 1-w=\sqrt{0.04}=0.2 \;\Rightarrow\; w=0.8, \nonumber\\
  \text{threshold } r:&\quad \frac{1}{1-r}\int_r^1\frac{(1-t)^2}{2}\,dt=\frac{(1-r)^2}{6}=0.02 \;\Rightarrow\; 1-r=\sqrt{0.12}=0.3464 \;\Rightarrow\; r=0.6536. \nonumber
\end{align}
At $v_1=r$, $g(r)=0.12/2=0.06>0.02$. The preemptive bid is
\begin{align}
  \bar p(r) &= v\cdot\Pr(\tilde v_2<v)+\int_v^r t\,dt + r\cdot\Pr(\tilde v_2>r) \nonumber\\
            &= (0.6)(0.6)+\frac{r^2-0.36}{2}+r(1-r) = 0.36+0.0336+0.2264 = 0.62. \nonumber
\end{align}
Types $v_1\in[0.6536,\,0.8)$ deter bidder 2 by pooling at $0.62$ but would face entry under observability.
\answer{(observable $v_1$)} Bidder 1 bids $v$; bidder 2 competes iff $v_1<w$, with $w>r$. Types in
$[r,w)$ are deterred in the signalling equilibrium but not under observability: preemption is more
frequent with private information, and the target receives a premium ($\bar p(r)>v$) only because of it.

\intuition Signalling is needed only to the extent that it changes the rival's action, and that
action is binary (enter or not). A binary decision can be steered by a binary message, so the
equilibrium pools types. Private information lets weak-but-above-$r$ types free ride on strong types'
deterrence.
\pitfalls
\begin{itemize}
  \item Expecting separation because types differ: separation requires each type's payoff to depend on the
        rival's \emph{exact} belief, which it does not here.
  \item Computing $r$ from $g(r)=c_2$ (the marginal type) instead of the average over $v_1\ge r$.
\end{itemize}
\bookans Matches the book's answer (p.~276): bidder 2 learns only $v\le v_1<r$ or $v_1\ge r$; with
observable $v_1$ it competes iff $v_1<w$ with $r<w$, so types in $[r,w)$ are deterred only in the
signalling equilibrium.
\refbib{Amaro de Matos \S8.2.1, pp.~206--207; Fishman (1988).}

% =============================================================================
\exsec{10. Jegadeesh--Chowdhry: no pooling equilibrium --- Amaro de Matos 8.1.5 \#8, p.~205}
%% BEGIN-STATEMENT 10
\stmtbegin
\textbf{Source:} Amaro de Matos, Ch. 8, Exercise 8.1.5 \#8, p.~205 (PDF 223) · Answer in book: yes,
pp.~274--275

8. In the discussion of a pooling equilibrium in the context of Jegadeesh and Chowdhry (1994),
assume that \(\alpha_m = 0.3\) and that there are two types of bidders with \(Z_L = 40\) and
\(Z_H = 80\). Assume further that 60\% are of the high-\(Z\) type and that \(P(\pi) = \ln(\pi)\).
Conditional on this information, show that there is no pooling equilibrium.

\paragraph{Context needed to solve it}

\emph{From Amaro de Matos §8.1.2 ``The Role of a Large Shareholder'', p.~196 (PDF 214) --- notation
shared by §8.1.4:}

{[}\ldots{]} Let \(\alpha\) be the proportion of the shares owned by L. Since he/she is a minority,
\(0 < \alpha < \frac{1}{2}\). Also, let \(Z\) denote the random increase in the discounted market
value generated by the substitution of an inefficient management. {[}\ldots{]} Under these
conditions, the bidder will move to make an offer if he/she can purchase \(\frac{1}{2} - \alpha\) of
the shares. The expected market value under the existing management being \(q\), the tender price
must be higher than \(q\), say, \(p = q + \pi\) {[}\ldots{]}.

\emph{From Amaro de Matos §8.1.4 ``The Optimal Size of α before a Takeover'', pp.~201--204 (PDF
219--222):}

Until now, the number of shares held by L has been considered as given. However, according to the
preceding reasoning, the proportion \(\alpha\) is likely to affect the probability that a takeover
is successful. Hence, it is a strategic variable, the optimal value of which can be established
before a takeover attempt. This problem was formalized by Jegadeesh and Chowdhry (1994), assuming
that the initial shareholdings is an endogenous variable to be determined so as to maximize the
value of the firm, an assumption that replaces Assumption 8.1.2.

It is assumed that the bidder buys a fraction \(\alpha\) of the shares in the open market at a
current market price of zero. The value of \(\alpha\) will be taken as a signal of his/her
motivation to take the control. In his/her attempt to acquire a controlling interest, the bidder can
purchase up to \(\alpha_{\max}\) shares in the open market, where \(\alpha_{\max} < 0.5\), and the
actual number of shares \(\alpha \in [0, \alpha_{\max}]\) that the bidder acquires in the open
market is determined endogenously.

The risk-neutral shareholder will not tender unless the bid is at least as large as the expected
improvements of the profits, where both the bidder's toehold and the bid are considered as a signal;
that is,

\[\pi \geq E(Z \mid \alpha, \pi).\]

This equation expresses the so-called \emph{free-rider condition}. The value of the bid for the
bidder is therefore

\[G(\pi, \alpha; Z) = P(\pi)[0.5Z - (0.5 - \alpha)\pi], \qquad (8.6)\]

and his/her problem is to maximize this value,

\[\max_{\pi, \alpha} G(\pi, \alpha; Z) \text{ subject to } \pi \geq E(Z \mid \alpha, \pi).\]

To obtain insightful results, two additional assumption are made:

\textbf{Assumption 8.1.4.} \emph{The probability of a successful bid is increasing and concave in
\(\pi\).}

{[}\ldots{]}

\textbf{Assumption 8.1.5.} \emph{The constraint imposed by the free-rider condition is binding for
all bidder types.}

Formally,

\[Z > \hat\pi(\alpha_{\max}, Z) > \hat\pi(\alpha, Z).\]

{[}\ldots{]} Their first result states that, in the absence of the free-rider constraint, the value
of bidding high is increasing in the bidder's type \(Z\). {[}\ldots{]} The second result says that,
in the absence of the free-rider constraint, the cost of choosing any level of initial \(\alpha\) is
identical for the bidder regardless of his/her value of \(Z\). {[}\ldots{]}

The type of equilibrium obtained in this bid is discussed in detail by Jegadeesh and Chowdhry. There
will be no pooling equilibrium---that is, an equilibrium in which \(\alpha\) has no value as a
signal and \(\hat\pi\) is equal for all the bidders.

To see why such an equilibrium would not hold, suppose there are two bidders, with different values
of \(Z\)---say, \(Z_L\) and \(Z_H\) with \(Z_L < Z_H\). Assume that the unconditional expected value
of \(Z\) is \(E(Z) = Z_M\), such that \(Z_L < Z_M < Z_H\).

Since in a pooling equilibrium the minimum value of the bid has to be the expected value of \(Z\) in
order to avoid the free-rider problem and \(\alpha\) plays no signaling role, every bidder will
choose to bid \(Z_M\) with \(\alpha_m\) in hand.

This equilibrium does not hold because the bidder with \(Z_L\) is bidding above his/her valuation.
His/her strategy to escape from that situation is to bid \(Z_L\) with a lower value of \(\alpha\)
(say, \(\alpha'\)) to be chosen such that the other bidder does not have an incentive to mimic
him/her. That is to say, \(\alpha'\) must satisfy

\[G(Z_L, \alpha'; Z_H) = G(Z_M, \alpha_m; Z_H).\]

It may be shown that in such a situation the low-valuation bidder is strictly better off, which is
to say

\[G(Z_L, \alpha'; Z_L) > G(Z_M, \alpha_m; Z_L)\]

and hence the pooling equilibrium does not survive. {[}\ldots{]}
\stmtend
%% END-STATEMENT

\asked With toehold $\alpha_m=0.3$, types $Z_L=40$ (40\%) and $Z_H=80$ (60\%), and success probability
$P(\pi)=\ln\pi$, show that the pooling equilibrium (both types hold $\alpha_m$ and bid $Z_M$) breaks down.

\approach A pooling equilibrium fails if some type can deviate to an action that the other type would
not mimic and that the deviating type strictly prefers. Here the low type deviates to a lower bid
$Z_L$ together with a smaller toehold $\alpha'$, chosen so that the high type is exactly indifferent
(the cheapest credible signal). This uses results~1 and 2 of Jegadeesh--Chowdhry: the cost of a lower
toehold is the same for both types, but a lower bid helps the low type more.

\notation
\begin{itemize}
  \item $G(\pi,\alpha;Z)=P(\pi)[0.5Z-(0.5-\alpha)\pi]$: bidder's value (8.6), bid premium $\pi$, toehold $\alpha$, type $Z$.
  \item $P(\pi)=\ln\pi$: increasing and concave (Assumption~8.1.4) and positive for $\pi>1$.
  \item $Z_M=\E[Z]$: pooled bid (free-rider condition binding: $\pi=\E[Z\mid\alpha,\pi]$).
\end{itemize}

\step{1}{Pooled bid} With no information in $\alpha$ or $\pi$, the free-rider condition requires
$\pi=\E[Z]$:
\begin{align}
  Z_M &= (0.4)(40)+(0.6)(80) = 16+48 = 64. \label{r9e10-ZM}
\end{align}

\step{2}{Values in the pool} With $\ln 64=4.15888$,
\begin{align}
  G(64,0.3;80) &= \ln(64)\,[0.5(80)-(0.5-0.3)(64)] \nonumber\\
               &= 4.15888\,[40-12.8] = 4.15888(27.2)=113.1216, \label{r9e10-GH}\\
  G(64,0.3;40) &= \ln(64)\,[0.5(40)-(0.2)(64)] \nonumber\\
               &= 4.15888\,[20-12.8] = 4.15888(7.2)=29.944. \label{r9e10-GL}
\end{align}
The low type pays $64>40$ for a gain worth 40 to it: it is subsidizing the high type.

\step{3}{The deviation that the high type will not mimic} The low type bids $\pi=Z_L=40$ (which
shareholders accept if they believe the bidder is the low type) with toehold $\alpha'$ such that the high
type is indifferent:
\begin{align}
  G(40,\alpha';80) &= G(64,0.3;80) \nonumber\\
  \ln(40)\,[40-(0.5-\alpha')40] &= 113.1216 \nonumber\\
  3.688879\,(20+40\alpha') &= 113.1216 \nonumber\\
  20+40\alpha' &= \frac{113.1216}{3.688879} = 30.6656 \nonumber\\
  \alpha' &= \frac{30.6656-20}{40} = 0.26664. \label{r9e10-ap}
\end{align}
Any $\alpha$ slightly below $\alpha'$ makes mimicking strictly unattractive for the high type.

\step{4}{The low type gains} Its value from the deviation is
\begin{align}
  G(40,\alpha';40) &= \ln(40)\,[0.5(40)-(0.5-0.26664)(40)] \nonumber\\
                   &= 3.688879\,[20-9.3344] = 3.688879(10.6656) = 39.344 > 29.944. \label{r9e10-dev}
\end{align}
\answer{(\#8)} The low type strictly prefers to bid 40 with a toehold $\alpha'=0.2666<0.3$ (value $39.34$
versus $29.94$ in the pool), and the high type does not imitate it (value $113.12$ either way); so the
pooling equilibrium does not survive. The toehold works as a signal: a \emph{smaller} toehold credibly
signals a lower valuation.

\step{5}{Why the signal separates} Result 2: $\partial G/\partial\alpha=\pi P(\pi)$ does not depend on $Z$, so
reducing the toehold from $0.3$ to $\alpha'$ at bid 40 costs both types the same,
$40\ln(40)(0.3-0.26664)=4.92$. Result 1: $\partial[G(64)-G(40)]/\partial Z=0.5[P(64)-P(40)]>0$, so
bidding low is worth more to the low type. Equal costs and unequal benefits sort the types.

\intuition In Jegadeesh--Chowdhry the bid premium \emph{rises} with the toehold: a bidder who buys
many shares before the offer reveals a high valuation, and must then bid high. A low-valuation bidder
deliberately buys fewer shares than it could, to be believed when it bids low.
\pitfalls
\begin{itemize}
  \item Choosing $\alpha'$ to make the \emph{low} type indifferent; the binding constraint is that the high type does not mimic.
  \item Using $Z_M$ as the low type's bid after deviating; the point of the deviation is the lower bid.
\end{itemize}
\bookans Matches the book's answer (pp.~274--275): $Z_M=64$, $G=113.1216$ and $29.944$,
$\alpha'=0.2667$; the book prints $G(40,0.2667;40)=39.337$, against $39.344$ here with $\alpha'$ unrounded
(rounding of $\alpha'$ and of the logarithms).
\refbib{Amaro de Matos \S8.1.4, pp.~201--204; Jegadeesh and Chowdhry (1994).}

% =============================================================================
\exsec{11. Takeover defenses: a flip-over pill whose new shares vote --- Tirole 11.1, p.~441}
%% BEGIN-STATEMENT 11
\stmtbegin
\textbf{Source:} Tirole, Ch. 11, Exercise 11.1, p.~441 (PDF 454) · Answer in book: no

\textbf{Exercise 11.1 (takeover defenses).} Extend the analysis of takeover defenses in Section
11.5.2 to the case in which the new shares created by the flip-over plan carry a voting right.

\paragraph{Context needed to solve it}

\emph{From Tirole §11.5 ``Positive Theory of Takeovers: Single-Bidder Case'', pp.~430--431 (PDF
443--444):}

{[}\ldots{]} In a first step we will assume that all shares carry equal voting rights and that the
raider needs a simple majority or, more generally, a fraction \(\kappa \in (0, 1)\) of the shares in
order to gain control, replace the incumbent management, and implement the new policy. That is, a
raider who purchases only a minority of shares or, more generally, a fraction less than \(\kappa\)
of the shares is on the same footing as any other investor, and neither delivers investor value
\(\hat v\) nor enjoys rent \(\hat w\) from control. {[}\ldots{]}

We say that the raider enhances value (to investors) if \(\hat v > v\). The case \(\hat v < v\)
corresponds to a ``value-decreasing raider.'' We focus on the case in which the value enhancement or
decrease is symmetric information.

\emph{(\(v\) is the investor value of the firm under incumbent management, \(\hat v\) under raider
management; §11.5.3, p.~434: ``Each share carries a cash-flow right equal to \(1/n\)th of the
investor payoff (\(v\) under incumbent management, \(\hat v\) under raider management).'')}

\emph{From Tirole §11.5.1 ``Value-Enhancing Raider: The Grossman--Hart Analysis'', pp.~431--432 (PDF
444--445):}

{[}\ldots{]} let us normalize to 1 the value added by the raider's taking over:

\[\hat v - v \equiv 1.\]

And let us redefine \(P\) as the premium over \(v\) offered by the raider. That is, the raider
offers price \(v + P\). The relevant range for \(P\) is \([0, 1]\). A negative-premium offer is
always rejected, while premia above 1 are accepted but wasteful for the raider.

We will assume that there is a continuum, of mass 1, of shareholders. The assumption of a continuum
finesses the issue of a shareholder's potentially being ``pivotal,'' i.e., affecting the outcome
through his tendering choice. A shareholder will therefore compare the takeover premium that is
offered by the raider to the expected value enhancement, taking the probability of takeover as
exogenous to his tendering choice and equal to its equilibrium value. {[}\ldots{]}

Consider an unrestricted, unconditional offer, and assume that the raider needs to acquire a
fraction \(\kappa \in (0, 1)\) of the shares to gain control. We claim that the probability of
takeover success must be equal to the premium:

\[\beta \equiv \Pr(\text{takeover success}) = P.\]

If this probability exceeded \(P\), then each shareholder would be better off holding on to his
share, since \(\beta\hat v + (1 - \beta)v > v + P\); and so the takeover would fail with probability
1, a contradiction. Similarly, for a probability of success smaller than \(P\), all would be better
off tendering and thus the takeover would succeed. So the equilibrium probability of takeover
success must equal the premium. The fraction of shares tendered must be exactly equal to \(\kappa\).
{[}\ldots{]} Leaving aside any private surplus \(\hat w\), the profit made by the raider on the
takeover attempt is

\[\pi = \kappa[\beta \cdot 1 - P] = 0.\]

That is, the raider is unable to derive any benefit from the value enhancement. {[}\ldots{]}

\emph{From Tirole §11.5.2 ``Positive Raider Surplus despite Free Riding'', pp.~432--434 (PDF
445--447):}

\textbf{11.5.2.1 Private Benefit from Control.} When the raider derives a private surplus \(\hat w\)
from control, then he gets to keep this surplus and optimally bids \(P = 1\). To see this, note that
the raider's profit when offering premium \(P\) is

\[\pi = \kappa[\beta - P] + \beta\hat w = P\hat w.\]

The raider strictly prefers to bid the maximum premium: \(P = 1\). {[}\ldots{]}

\textbf{11.5.2.2 Toehold.} {[}\ldots{]} Suppose that the raider already owns a fraction
\(\theta < \kappa\) of the shares when making a tender offer. Assuming, again, that \(\hat w = 0\),
the raider's profit for premium \(P\) is

\[\pi(P) = (\kappa - \theta)(\beta - P) + \theta\beta,\]

where \(\beta\) is, as earlier, the probability of takeover success and must in equilibrium be equal
to \(P\). Hence, \(\pi(P) = \theta P\). The optimal bid is then \(P = 1\), yielding profit
\(\pi = \theta\). {[}\ldots{]}

\textbf{11.5.2.4 Takeover Defenses.} Takeover defenses come in many guises, and, except for the
common feature that they make it harder for a raider to acquire a firm, are hard to summarize
concisely. Let us illustrate their role in the case of poison pills, more specifically in the most
common form of a ``flip-over plan'' under which the holders of shares are entitled to purchase new
shares at a substantial discount after a hostile takeover.

For computational simplicity, let us assume a simple majority rule (\(\kappa = \frac{1}{2}\)) and
that the new shares carry no voting rights.\textsuperscript{26} In the case of takeover success, the
50\% of shares kept by the initial shareholders are worth \(\hat v + \Delta\) (with \(\Delta > 0\)),
while the 50\% acquired by the raider are worth \(\hat v - \Delta\) to him due to the
dilution.\textsuperscript{27} Letting, as before, \(\beta\) denote the probability of success and
\(P\) the premium over \(v\), shareholders are indifferent between tendering and keeping their
shares if and only if

\[\beta(\hat v + \Delta) + (1 - \beta)v = v + P \quad\text{or}\quad \beta = \frac{P}{1 + \Delta}.\]

The raider's profit is then

\[\pi = \beta\hat w + \tfrac{1}{2}[\beta(\hat v - \Delta) + (1 - \beta)v - (v + P)] = \beta(\hat w - \Delta).\]

Assuming that \(\hat w > \Delta\) (otherwise the raider makes no offer), it is optimal for the
raider to succeed for certain (\(\beta = 1\)) by choosing \(P = 1 + \Delta\).

The poison pill further raises the purchase price. In contrast with the dilution of initial
shareholders by the raider considered in Grossman and Hart, poison pills allow a dilution of the
raider by initial shareholders. Poison pills thereby allow the firm to adjust the purchase price
paid by the raider. Suppose, for example, that the raider's benefit from control, \(\hat w\), is
known (the distribution \(H\) is a spike at \(\hat w\)).\textsuperscript{28} In the absence of a
poison pill, \(P = 1\) and the raider's surplus is \(\hat w\). The optimal poison pill then yields
dilution \(\Delta^* = \hat w\).

\textsuperscript{26} If the new shares carry a voting right, the analysis is basically unchanged.

\textsuperscript{27} For example, suppose that the holders of untendered shares are entitled to one
extra share per share for free. Then the raider has only one-third of the cash rights. Then
\(\Delta = \frac{1}{3}\hat v\).

\textsuperscript{28} And \(v + w < \hat v + \hat w\).
\stmtend
%% END-STATEMENT

\asked Redo Tirole's poison-pill analysis (\S11.5.2.4) when the shares issued to non-tendering
shareholders after a hostile takeover carry votes: what fraction the raider must buy, the
shareholders' indifference condition, the raider's premium and profit, and the optimal pill.

\approach Free-riding logic (Grossman--Hart, Tirole \S11.5.1): with a continuum of shareholders the
success probability $\beta$ makes them indifferent between tendering and holding. The pill is a
transfer from the raider to the non-tendering holders after success. Voting new shares add one
element: they dilute the raider's \emph{votes}, so the raider needs a larger fraction of the old shares
to keep control after the pill is triggered. Tirole's footnote 26 says the analysis is ``basically
unchanged''; we verify that and quantify the change.

\notation
\begin{itemize}
  \item $v$, $\hat v$: investor value per old share under incumbent and raider, $\hat v-v\equiv 1$; $\hat w$: raider's private benefit.
  \item $m$: new shares received free by each untendered old share after a successful takeover (Tirole's example: $m=1$).
  \item $\kappa'$: fraction of old shares the raider must buy; $\Delta_u$: gain per untendered old share; $\Delta_r$: dilution per raider share.
  \item $P$: premium over $v$; $\beta$: probability of success.
\end{itemize}

\step{1}{Control threshold} After success the raider holds $\kappa'$ old shares and the others hold
$(1-\kappa')(1+m)$ shares, all voting. He keeps a majority iff
\begin{align}
  \frac{\kappa'}{\kappa'+(1+m)(1-\kappa')} &\ge \frac12 \nonumber\\
  2\kappa' &\ge \kappa'+(1+m)(1-\kappa') \nonumber\\
  \kappa'(2+m) &\ge 1+m \nonumber\\
  \kappa' &\ge \frac{1+m}{2+m}. \label{r9e11-k}
\end{align}
With non-voting new shares the threshold is $\kappa=\tfrac12$; with $m=1$ it becomes $\kappa'=2/3$.

\step{2}{Cash-flow split at the threshold} Total cash-flow shares are $\kappa'+(1+m)(1-\kappa')$, which
at $\kappa'=(1+m)/(2+m)$ equals $2(1+m)/(2+m)$. A raider share is worth
$\hat v(2+m)/[2(1+m)]$ and an untendered old share (with its $m$ new ones) $(1+m)\hat v(2+m)/[2(1+m)]=\hat v(2+m)/2$:
\begin{align}
  \Delta_u &= \frac{(2+m)\hat v}{2}-\hat v = \frac{m\hat v}{2}, &
  \Delta_r &= \hat v-\frac{(2+m)\hat v}{2(1+m)} = \frac{m\hat v}{2(1+m)}. \label{r9e11-D}
\end{align}
The pill is a pure transfer: $\kappa'\Delta_r=\frac{1+m}{2+m}\cdot\frac{m\hat v}{2(1+m)}=\frac{m\hat v}{2(2+m)}=(1-\kappa')\Delta_u$, so
\begin{align}
  \kappa'(\Delta_r+\Delta_u) &= (1-\kappa')\Delta_u+\kappa'\Delta_u = \Delta_u. \label{r9e11-cons}
\end{align}

\step{3}{Tendering indifference} A holder compares $v+P$ with keeping a share worth $\hat v+\Delta_u$ if the
raid succeeds and $v$ otherwise:
\begin{align}
  \beta(\hat v+\Delta_u)+(1-\beta)v &= v+P \nonumber\\
  \beta(\hat v-v+\Delta_u) &= P \nonumber\\
  \beta &= \frac{P}{1+\Delta_u}. \label{r9e11-beta}
\end{align}

\step{4}{Raider's profit} He buys $\kappa'$ shares at $v+P$, which are worth $\hat v-\Delta_r$ to him on
success and $v$ on failure:
\begin{align}
  \pi &= \beta\hat w+\kappa'[\beta(\hat v-\Delta_r)+(1-\beta)v-(v+P)] \nonumber\\
      &= \beta\hat w+\kappa'[\beta(1-\Delta_r)-P] \nonumber\\
      &= \beta\hat w+\kappa'[\beta(1-\Delta_r)-\beta(1+\Delta_u)] \quad\text{(using \eqref{r9e11-beta})} \nonumber\\
      &= \beta\hat w-\beta\kappa'(\Delta_r+\Delta_u) \nonumber\\
      &= \beta(\hat w-\Delta_u) \quad\text{(using \eqref{r9e11-cons})}. \label{r9e11-pi}
\end{align}
This is Tirole's formula with $\Delta$ replaced by $\Delta_u$. If $\hat w>\Delta_u$ the raider sets $\beta=1$,
i.e.\ $P=1+\Delta_u$, and earns $\hat w-\Delta_u$; otherwise he does not bid.

\step{5}{Numbers} $v=2$, $\hat v=3$, $\hat w=2$, $m=1$. Voting: $\kappa'=2/3$, $\Delta_u=1(3)/2=1.5$,
$\Delta_r=3/4=0.75$, check $\tfrac23(0.75)=\tfrac13(1.5)=0.5$; premium $P=1+1.5=2.5$ and profit
$2-1.5=0.5$. Non-voting (Tirole's footnote 27): $\kappa=\tfrac12$, $\Delta=\hat v\,m/(2+m)=3/3=1$, $P=2$,
profit $2-1=1$.

\step{6}{Optimal pill} With $\hat w$ known, the target sets $\Delta_u^*=\hat w$, extracting the whole private
benefit ($P=1+\hat w$, $\pi=0$), exactly as with non-voting shares. Only the instrument changes: voting new
shares deliver $\Delta_u=m\hat v/2$ per untendered share versus $m\hat v/(2+m)$, so fewer new shares are needed,
$m=2\hat w/\hat v=4/3$ (with $\kappa'=0.7$) instead of $m=2\hat w/(\hat v-\hat w)=4$ in the example.
\answer{(11.1)} With voting new shares the raider must buy $\kappa'=(1+m)/(2+m)>\tfrac12$ of the shares;
the equilibrium is otherwise unchanged: $\beta=P/(1+\Delta_u)$, $\pi=\beta(\hat w-\Delta_u)$, bid
$P=1+\Delta_u$ when $\hat w>\Delta_u$, optimal pill $\Delta_u^*=\hat w$. A given number of new shares
is a stronger defense (larger $\Delta_u$), and the pill also raises the number of shares that must be tendered.

\intuition Whatever the pill does to votes, it is a transfer from the raider to the holdouts on
success. Free riding prices that transfer into the premium, so the raider pays it on every share he
buys; net of the transfer, he keeps only his private benefit minus the holdouts' gain. Voting rights
change how many shares he must buy, not that accounting.
\pitfalls
\begin{itemize}
  \item Keeping $\kappa=\tfrac12$: with voting new shares, a bare majority of old shares loses control once the pill is triggered.
  \item Writing the raider's loss per share as $\Delta_u$ (the holdouts' gain); only at $\kappa=\tfrac12$ are the two equal.
\end{itemize}
\bookans No answer in the book; Tirole's footnote 26 (``basically unchanged'') is consistent with the result.
\refbib{Tirole \S11.5.1--11.5.2, pp.~431--434.}

% =============================================================================
\exsec{12. Hansen: every cash offer is dominated by a stock offer --- Amaro de Matos 8.2.6 \#4, p.~223}
%% BEGIN-STATEMENT 12
\stmtbegin
\textbf{Source:} Amaro de Matos, Ch. 8, Exercise 8.2.6 \#4, p.~223 (PDF 241) · Answer in book: no
(p.~277 only cites page 80 of Hansen 1987)

4. Given asymmetric information in the model of Hansen (1987), show that for every feasible cash
offer there exists a stock offer that is preferred by the acquirer.

\paragraph{Context needed to solve it}

\emph{From Amaro de Matos §8.2.2 ``Choosing the Means of Payment'', pp.~207--210 (PDF 225--228):}

In a model developed by Hansen (1987), payment is made either in cash or in stock of the merged
corporation. The basis of the mechanism that makes management and bidder prefer one or the other
(either cash or stock) is asymmetry of information.

The acquirer owns a firm of value \(x\), which is known to both the acquirer and the target's
management. The target firm is worth \(v\) to itself, but has a value \(w(v) \geq v\) to the
acquirer. The value of a merger will therefore be \(x + v\). The asymmetry enters assuming that the
bidder does not know the value of the firm \(v\), but the current management does. Then, the target
will never accept an offer below that value in a cash trade. This creates an adverse selection
problem, against which the acquirer can protect himself/herself, defining his/her optimal strategy
not on unconditional expected profit, but on expected profit conditional on the offer being
accepted.

The acquirer takes \(v\) as a random variable with support on the closed interval \([v_l, v_h]\).
For a cash offer \(C\), the expected profit of the raider is an expectation conditional on the
acceptance of the offer---that is to say, conditional on \(v_l < v < C\). If \(f(v)\) denotes the
density of probability of \(v\),

\[E(\Pi_r \mid C) = \int_{v_l}^{C} [w(v) - C] f(v)\, dv.\]

Denoting by \(F(v)\) the distribution of \(v\) and by \(E(w \mid C)\) the expected value of the
target to the acquirer conditional on sale,

\[E(w \mid C) = \frac{1}{F(C)} \int_{v_l}^{C} w(v) f(v)\, dv,\]

the expected conditional wealth gain can be rewritten as

\[E(\Pi_r \mid C) = F(C)[E(w \mid C) - C].\]

{[}\ldots{]} But if the bidder offers stock of the merged firm (say, a proportion \(\beta\)), trade
may still occur. From the point of view of the acquirer, such a trade is in the interest of the
target with certainty if the final value held by the owners of the target is greater than or equal
to \(v\) for whatever current value \(v \in [v_l, v_h]\). It will be in the interest of the target
with probability \(F(v^*)\) if the same inequality holds for any \(v \in [v_l, v^*]\). If the
management of the target knows the function \(w\), this means that with probability \(F(v^*)\) trade
is good for the target if

\[\beta[x + w(v)] \geq v \quad \forall v \in [v_l, v^*]\]

or

\[\beta \geq \frac{v}{x + w(v)} \quad \forall v \in [v_l, v^*], \qquad (8.11)\]

which is equivalent to

\[\beta \geq \max_v \frac{v}{x + w(v)}.\]

Making Hansen's additional assumption that

\[\frac{dw}{dv} < \frac{x + w(v)}{v}\]

the preceding condition becomes simply

\[\beta[x + w(v^*)] \geq v^*.\]

More easily, the acquirer will be interested only if the value held by him/her after the transaction
is greater than the current value. In other words,

\[(1 - \beta)[x + E(w \mid v^*)] \geq x. \qquad (8.12)\]

{[}\ldots{]} In a stock offer the expected wealth of the acquirer will be

\[F(v^*)(1 - \beta)[E(w \mid v^*) + x] + [1 - F(v^*)]x. \qquad (8.13)\]

{[}\ldots{]} The conditions under which cash or stock trades with a given probability are possible
are, therefore, established. Of course, it is possible that they are simultaneous. In such a case,
for each feasible cash offer there will be a stock offer with the same probability of success that
will be preferred by the acquirer. To see this, notice that the expected wealth of the acquirer
after a cash offer \(C\) equals

\[F(C)[E(w \mid C) - C] + x. \qquad (8.14)\]

Constructing a stock offer with \(\beta\) given by

\[\beta^* = \frac{C}{x + w(C)},\]

it is known from (8.11) that it will have the same probability of success \(F(C)\) as a cash offer
of value \(C\). The acquirer's expected wealth in this case will be

\[F(C)(1 - \beta^*)[E(w \mid C) + x] + [1 - F(C)]x.\]

It is a simple algebraic exercise to check that this expression is greater than the one in (8.14).
Therefore, in the situation described up to now in this model, no cash offers will occur, and only
stock trade is possible.
\stmtend
%% END-STATEMENT

\asked For any cash offer $C$ that the acquirer would make (non-negative expected gain), construct a
stock offer that is accepted by exactly the same target types and leaves the acquirer with a higher
expected wealth.

\approach The book's construction: $\beta^*=C/[x+w(C)]$. Two steps: (i) show that under Hansen's
assumption $dw/dv<(x+w)/v$ the target types who accept $\beta^*$ are exactly those who accept $C$,
namely $v\le C$; (ii) subtract the two expected wealths and sign the difference. The economics: stock
pays the target in a currency whose value moves with $w(v)$, which shifts part of the lemons cost back to
the target.

\notation
\begin{itemize}
  \item $v\in[v_l,v_h]$, distribution $F$, density $f$; $w(v)\ge v$, increasing; $x$: acquirer's value (known).
  \item $C$: feasible cash offer, $E(\Pi_r\mid C)=F(C)[E(w\mid C)-C]\ge 0$, with $E(w\mid C)=\E[w(v)\mid v\le C]$.
  \item $\beta^*=C/[x+w(C)]$; $W^{\text{cash}}$, $W^{\text{stock}}$: acquirer's expected wealth.
\end{itemize}

\step{1}{The acceptance rule is the same} Type $v$ accepts $\beta^*$ iff $\beta^*[x+w(v)]\ge v$, i.e.\ iff
$h(v)\equiv v/[x+w(v)]\le\beta^*=h(C)$. Differentiating $h$ (quotient rule),
\begin{align}
  h'(v) &= \frac{[x+w(v)]-v\,w'(v)}{[x+w(v)]^2} > 0 \iff w'(v) < \frac{x+w(v)}{v}, \label{r9e12-h}
\end{align}
which is Hansen's assumption. So $h$ is strictly increasing and $h(v)\le h(C)\iff v\le C$: the stock offer
is accepted with probability $F(C)$, by the same types as the cash offer.

\step{2}{The two expected wealths} From (8.14) and the stock formula,
\begin{align}
  W^{\text{cash}} &= F(C)[E(w\mid C)-C]+x, \label{r9e12-Wc}\\
  W^{\text{stock}} &= F(C)(1-\beta^*)[E(w\mid C)+x]+[1-F(C)]x. \label{r9e12-Ws}
\end{align}

\step{3}{The difference} Expand \eqref{r9e12-Ws} and subtract \eqref{r9e12-Wc}:
\begin{align}
  W^{\text{stock}}-W^{\text{cash}} &= F(C)\{(1-\beta^*)[E(w\mid C)+x]-x\}+x-F(C)[E(w\mid C)-C]-x \nonumber\\
  &= F(C)\{E(w\mid C)+x-\beta^*[E(w\mid C)+x]-x-E(w\mid C)+C\} \nonumber\\
  &= F(C)\{C-\beta^*[E(w\mid C)+x]\}. \nonumber
\end{align}
Substitute $C=\beta^*[x+w(C)]$ (definition of $\beta^*$):
\begin{align}
  W^{\text{stock}}-W^{\text{cash}} &= F(C)\{\beta^*[x+w(C)]-\beta^*[E(w\mid C)+x]\} \nonumber\\
  &= F(C)\,\beta^*\,[w(C)-E(w\mid C)]. \label{r9e12-diff}
\end{align}

\step{4}{Sign} $E(w\mid C)$ averages $w$ over $v\le C$, and $w$ is increasing, so $E(w\mid C)\le w(C)$ with
strict inequality when $w$ is strictly increasing and $F$ has mass below $C$. Since $F(C)>0$ and
$\beta^*>0$ for an offer that is ever accepted, \eqref{r9e12-diff} is positive. The stock offer also
satisfies the acquirer's participation constraint, because its gain exceeds the cash gain, which is $\ge 0$.

\step{5}{Illustration} $v\sim U[0,100]$, $w(v)=1.2v+30$, $x=200$ (Hansen's condition holds:
$1.2<(230+1.2v)/v$). For $C=60$: $F=0.6$, $E(w\mid C)=1.2(30)+30=66$,
\begin{align}
  \text{cash gain} &= 0.6(66-60)=3.6, \qquad \beta^*=\frac{60}{200+102}=0.19868, \nonumber\\
  \text{stock gain} &= 3.6+0.6(0.19868)(102-66)=3.6+4.29=7.89. \nonumber
\end{align}
The optimal cash offer solves $w(C)-C-F(C)/f(C)=1.2C+30-C-C=0$, so $C^*=37.5$ with gain
$0.375(52.5-37.5)=5.625$; the matching stock offer does better again by $0.375\,\beta^*(75-52.5)$.
\answer{(\#4)} For any feasible $C$, the stock offer $\beta^*=C/[x+w(C)]$ is accepted by the same types
($v\le C$) and raises the acquirer's expected wealth by $F(C)\beta^*[w(C)-E(w\mid C)]\ge 0$. Cash offers are
(weakly) dominated; strictly so when $w$ is strictly increasing.

\intuition With cash the marginal accepting type $v=C$ is paid $C$, and so is every lower type.
With stock, the payment $\beta^*[x+w(v)]$ is lower for lower types, because the shares they receive are
worth less when the asset they bring is worth less. The acquirer saves $\beta^*[w(C)-w(v)]$ on each
inframarginal type, which is \eqref{r9e12-diff} on average.
\pitfalls
\begin{itemize}
  \item Matching the stock offer by value at the mean type instead of at the marginal type $v=C$; then the acceptance sets differ.
  \item Forgetting Hansen's monotonicity assumption: without it $v/[x+w(v)]$ may be non-monotone and the acceptance set of a stock offer is not an interval.
\end{itemize}
\bookans The book only refers to page 80 of Hansen (1987) (p.~277); the proof above completes the
``simple algebraic exercise'' stated on p.~210.
\refbib{Amaro de Matos \S8.2.2, pp.~208--210; Hansen (1987).}

% =============================================================================
\exsec{13. Fishman (1989): face value needed to signal, cash versus debt --- Amaro de Matos 8.2.6 \#8, p.~223}
%% BEGIN-STATEMENT 13
\stmtbegin
\textbf{Source:} Amaro de Matos, Ch. 8, Exercise 8.2.6 \#8, p.~223 (PDF 241) · Answer in book: yes,
pp.~277--278

8. Show that, in the model of Fishman (1989), the face value needed to signal that \emph{v}₁ is at
or above any given level is lower with cash compared to debt.

\paragraph{Context needed to solve it}

\emph{Fishman (1988) base model --- from Amaro de Matos §8.2.1, pp.~206--207 (PDF 224--225):}

{[}\ldots{]} there are two bidders (say, bidder 1 and bidder 2), who make successive bids until only
one remains. The remaining bid is accepted by the target, provided that its value is greater than or
equal to the value \(v\), which is assumed to be reasonable by the target. Otherwise, there is no
takeover. Here, \(v\) is taken as the current market value of the target, assuming that there is no
takeover.

Let bidder 1 be the first to make an offer. The costs \(c_i\) for \(i = 1, 2\) of the bid are the
costs to the bidders of private observation of their own valuations \(v_i\), \(i = 1, 2\).
Furthermore, it is assumed that the valuations are random variables such that \(E(\tilde v_i) < v\).
{[}\ldots{]} If \(p\) is the initial offer made by bidder 1, it must be such that
\(v \leq p \leq v_1\) and at the same time makes competition unattractive to bidder 2. In other
words, the expected payoff from competing to bidder 2 should be nonpositive.

The payoffs can be expressed simply in terms of a variable \(d\) that takes the value 0 if bidder 2
does not compete and takes 1 otherwise. For any two random realizations of \(v_1\) and \(v_2\) the
payoff for bidder 1 is

\[\Pi_1(v_1, v_2, p, d) = v_1 - (1 - d)p - d \min[\max(p, v_2), v_1] - c_1 \qquad (8.9)\]

and for bidder 2 is

\[\Pi_2(v_1, v_2, p, d) = d[(v_2 - \min(v_2, v_1) - c_2]. \qquad (8.10)\]

{[}\ldots{]} both of them know the densities of probabilities \(f_i(v_i)\), which are assumed to be
strictly positive in a closed interval \([l, h]\) and zero elsewhere. {[}\ldots{]} Denoting this
threshold by \(r\), it must be such that the expected payoff of bidder 2 conditional on
\(v_1 \geq r\) is zero {[}\ldots{]}.

\emph{From Amaro de Matos §8.2.3 ``Cash as a Preemptive Instrument with Many Bidders'', pp.~212--214
(PDF 230--232):}

In the paper of Hansen, there is only one acquirer for the target. Fishman (1989) developed a model
for choice of means of payment that allows more potential acquirers. In this model it is possible to
see how the presence of competing bidders may affect the choice of the means of payment differently.

His idea was to extend his model of preempting bids discussed before, making the choice of the means
of payment a reinforcing preemptive device.

The model is basically Fishman's earlier model, but the profit of the bidders will depend on an
additional variable \(\theta \in \{1, 0\}\) describing the chosen means of payment (1 for cash and 0
for issuing debt). In the event of a debt offer, \(p\) will denote its face value.

Since the value of a debt offer is state contingent, it is assumed that the target costlessly
observes a private signal \(s\) that conveys information on both bidders' valuations. It is supposed
that \(s\) may assume only two possible values, \(s^+\) and \(s^-\), referring respectively to a
``good'' and a ``bad'' state of the world, with probabilities \(\gamma\) and \(1 - \gamma\). The
valuations of the bidders will also depend on the state of the world that they infer. In other
words, it is as if the density functions \(f_i\) depended on \(s\). The value of any variable of
interest will be labeled with a ``+'' if it is conditional on \(s = s^+\) and with a ``−'' if
conditional on \(s = s^-\). For instance, it is assumed that \(v_i^+ > v_i^-\). That is to say, for
any bidder the valuation is higher in the good state of the world. Hence, the assumption that any
bidder 1's offer face value is between \(v\) and \(v_1\) may read

\[v \leq p \leq v_1^+.\]

In particular, a debt offer with face value \(p\) is worth \(\min\{v_i^+, p\} \geq v\) if
\(s = s^+\) and \(\min\{v_i^-, p\} < v\) if \(s = s^-\). Thus, the target will accept the offer if
\(s = s^+\) and will reject it if \(s = s^-\).

The main difference between a debt and a cash offer is that the value of the cash offer does not
depend on the state of the world. Thus, cash offers cannot induce the target to make an efficient
acceptance/rejection decision.

The profit itself will depend, directly and indirectly, on the state of the world. Introducing a
variable \(\delta\) that equals 1 when \(s = s^+\) and 0 when \(s = s^-\), the profits in (8.9) and
(8.10) are replaced by

\[\Pi_1(v_1, v_2, \delta, p, \theta, d) = [\theta + (1 - \theta)\delta]v_1 - [\theta(1 - \delta) + \delta(1 - d)]p - d\delta \min[\max(p, v_2), v_1] - c_1\]

\[\Pi_2(v_1, v_2, \delta, p, \theta, d) = d\{\delta[v_2 - \min(v_2, v_1)] - c_2\} \qquad (8.15)\]

for bidder 1 and bidder 2, respectively.

As before, bidder 1 tries to find the threshold \(r\) for the firm value above which bidder 2 will
have zero profit. In other words, \(r\) is implicitly defined through

\[\frac{\int_r^h E_{v_2}\Pi_2(v_1, \tilde v_2, \delta, p, \theta, 1) f_1(v_1)\, dv_1}{\int_r^h f_1(v_1)\, dv_1} = 0.\]

For any bid \(v'\) higher than \(r\), the expected payoff of bidder 2 will be negative. On the other
hand, bidder 1 will be able to preempt only if \(v_1 > r\). Therefore, the expected gain of
preempting at an offer price \(p'\), other things constant, is

\[E\Pi_1(v_1, \tilde v_2, \delta, p', \theta, 0) - E\Pi_1(v_1, \tilde v_2, \delta, p', \theta, 1) = \gamma\{E_{v_2}\min[\max(p', \tilde v_2^+), v_1^+] - p'\}.\]

This expression is clearly increasing in \(v_1\). Hence, high-valuing first bidders stand to lose
more in an auction and thus have a greater incentive to deter competition.

On the other hand, the expected gain of making a debt offer compared to a cash offer at face value
\(p'\), other things constant, is

\[E\Pi_1(v_1, \tilde v_2, \delta, p', 0, d) - E\Pi_1(v_1, \tilde v_2, \delta, p', 1, d) = (1 - \gamma)[p' - v_1^-].\]

This gain is decreasing in \(v_1\). From these two results it appears that first bidders will tend
to preempt with cash whenever they have a high valuation for the target, and will possibly make a
debt offer if they do not intend to preempt. {[}\ldots{]}

Now, let us cross these two pieces of information. Let \(p_d(v_1)\) denote the highest face value of
an offer that bidder 1 is willing to make in order to preempt. It has to be such that the preceding
gain of preempting is zero. As we notice that the gain is decreasing in \(p'\), it follows that

\[p_d(v_1) = E\min[\max(v, \tilde v_2^+), v_1^+].\]

As for a cash offer, its maximum value \(p_c(v_1)\) is defined by a zero expected profit from both
preempting and paying in cash, which is to say, \(p_c(v_1)\) must satisfy

\[\gamma\{E_{v_2}\min[\max(p_c(v_1), \tilde v_2^+), v_1^+] - p_c(v_1)\} = (1 - \gamma)[p_c(v_1) - v_1^-]\]

or

\[p_c(v_1) = \gamma E_{v_2}\min[\max(v, \tilde v_2^+), v_1^+] + (1 - \gamma)v_1^-.\]

It is possible to show that there exists a perfect sequential equilibrium such that the optimal
strategy of bidder 1 is to make a debt offer of value \(v\) if he/she cannot preempt (which is to
say, if \(v \leq v_1 \leq r\)) or to make a cash offer with value \(p_c(r)\) if preemption is
possible (in other words, if \(v_1 > r\)).

The optimal strategy of bidder 2, on the other hand, will be to compete if bidder 1 made a debt
offer with face value lower than \(p_d(r)\) or if he/she made a cash offer of value lower than
\(p_c(r)\). If bidder 1's cash offer equals (or is greater than) \(p_c(r)\), bidder 2 does not
compete.
\stmtend
%% END-STATEMENT

\asked Show that, to signal that $v_1$ is at least a given level, bidder 1 needs a smaller face value
if it pays in cash than if it pays with debt: $p_c(v_1)<p_d(v_1)$.

\approach A signal works if lower types are unwilling to send it; the face value needed to signal
$v_1\ge r$ is the most that type $r$ is willing to offer, $p_c(r)$ with cash and $p_d(r)$ with debt
(above it type $r$ prefers not to preempt). Compare the two closed forms. The economics: cash is a
\emph{costlier} signal per unit of face value, because it is paid even in the bad state, when the
target is worth $v_1^-<p$ to the bidder; debt is refused in the bad state and costs nothing there.

\notation
\begin{itemize}
  \item $\gamma\in(0,1)$: probability of the good signal $s^+$; $v_1^\pm$, $\tilde v_2^\pm$: valuations in each state, $v_i^+>v_i^-$.
  \item $p_d(v_1)=\E\min[\max(v,\tilde v_2^+),v_1^+]$; $p_c(v_1)=\gamma\,p_d(v_1)+(1-\gamma)v_1^-$.
  \item Model facts: $v\le p\le v_1^+$ (so $v_1^+\ge v$) and $\min\{v_1^-,p\}<v$ with $p\ge v$ (so $v_1^-<v$).
\end{itemize}

\step{1}{Write the difference} Subtract $p_d$ from $p_c$:
\begin{align}
  p_c(v_1)-p_d(v_1) &= \gamma p_d(v_1)+(1-\gamma)v_1^- - p_d(v_1) \nonumber\\
  &= (1-\gamma)v_1^- - (1-\gamma)p_d(v_1) \nonumber\\
  &= (1-\gamma)\,[v_1^- - p_d(v_1)]. \label{r9e13-diff}
\end{align}

\step{2}{Bound $p_d$ from below} For every realization, $\max(v,\tilde v_2^+)\ge v$ and $v_1^+\ge v$, so
$\min[\max(v,\tilde v_2^+),v_1^+]\ge v$. Taking expectations, $p_d(v_1)\ge v$.

\step{3}{Bound $v_1^-$ from above} In the bad state the debt offer (face $p\ge v$) is worth
$\min\{v_1^-,p\}<v$. If $v_1^-\ge v$ then $\min\{v_1^-,p\}\ge v$, a contradiction; so $v_1^-<v$.

\step{4}{Sign} Combining, $v_1^-<v\le p_d(v_1)$, so $v_1^--p_d(v_1)<0$; with $1-\gamma>0$,
\eqref{r9e13-diff} is negative:
\begin{align}
  p_c(v_1) &< p_d(v_1). \label{r9e13-res}
\end{align}
For $\gamma=1$ (the target's signal is uninformative about the bad state), $p_c=p_d$.

\step{5}{Illustration} $\gamma=0.6$, $v=0.6$, $\tilde v_2^+\sim U[0,1]$, $v_1^+=0.9$, $v_1^-=0.4$:
\begin{align}
  p_d &= (0.6)(0.6)+\frac{0.9^2-0.6^2}{2}+(0.9)(0.1) = 0.36+0.225+0.09=0.675, \nonumber\\
  p_c &= (0.6)(0.675)+(0.4)(0.4) = 0.405+0.16=0.565, \nonumber
\end{align}
and $p_c-p_d=(0.4)(0.4-0.675)=-0.11$.
\answer{(\#8)} $p_c(v_1)-p_d(v_1)=(1-\gamma)[v_1^--p_d(v_1)]<0$ because $p_d\ge v>v_1^-$: signalling a
given valuation needs a lower face value with cash than with debt.

\intuition A signal is credible when it is expensive for those who should not send it. Paying cash commits
the bidder to overpay in the bad state (pay $p$ for an asset worth $v_1^-$), a cost that a low type fears
more; so a smaller cash offer is as informative as a larger debt offer. This is why, in Fishman's model,
preemptive (high-valuation) bidders pay cash, the opposite of Hansen's link between cash and an
undervalued acquirer.
\pitfalls
\begin{itemize}
  \item Comparing the \emph{values} of the offers to the target instead of their face values.
  \item Forgetting that debt is worthless to the bidder's cost in the bad state because the target rejects it.
\end{itemize}
\bookans Matches the book's answer (pp.~277--278): the inequality reduces to $(\gamma-1)[p_d-v_1^-]<0$,
equality when $\gamma=1$.
\refbib{Amaro de Matos \S8.2.3, pp.~212--214; Fishman (1989).}

% =============================================================================
\exsec{14. Shleifer--Vishny with an arbitrary resale price --- Tirole 14.2, p.~515}
%% BEGIN-STATEMENT 14
\stmtbegin
\textbf{Source:} Tirole, Ch. 14, Exercise 14.2, p.~515 (PDF 528) · Answer in book: yes, p.~618

\textbf{Exercise 14.2 (alternative distributions of bargaining power in the Shleifer--Vishny
model).} Perform the analysis of Section 14.2.2 for an arbitrary unit price \(P \in [0, \rho_0]\) of
resale of a distressed firm's assets to a productive one. (Assume that bargaining occurs between the
two firms' investors, and that the acquiring firm's investors then redesign their managerial
incentives. Thus the per-unit surplus \(\rho_0 - P\) goes to the acquiring firm's investors.)

\paragraph{Context needed to solve it}

\emph{From Tirole §14.2.2 ``Industry-Wide Shocks and Distress Sales: The Shleifer--Vishny Model'',
pp.~500--503 (PDF 513--516); equations checked against the PDF images:}

This subsection, building on Shleifer and Vishny (1992), endogenizes the resale price in the
two-firm, continuous-investment version of the model of Section 4.3.1. Let us restate the key
ingredients of this model.

\emph{Investment and redeployability.} There are two firms in the industry. The ``industry'' is here
defined as a group of symmetric firms using the same equipment/assets. For simplicity, we assume
that the two firms do not compete in the same product market (see below for a discussion of this
hypothesis). Each firm is run by an entrepreneur, who has initial cash \(A_i\).

Initially, firm \(i\) invests \(I_i\), and therefore borrows \(I_i - A_i\) from some \emph{ex ante}
competitive lender (``lender \(i\)''). Then, there is a costless ``learning period.'' At the end of
the learning period, each firm learns whether it is ``productive'' (which has probability \(x\)), or
``unproductive,'' i.e., ``in distress'' (which has probability \(1 - x\)). Being unproductive means
that the firm will always be unsuccessful regardless of whether the entrepreneur behaves. For
example, there may be no demand for the firm's output. Its assets are then useless if left in place.
A productive firm is described as in the variable-investment model of Section 3.4 (which,
incidentally, corresponds to the case \(x = 1\)).

If both firms are productive, each manages its initial investment. Firm \(i\)'s profit is either 0
or \(RI_i\). Borrower \(i\)'s private benefit is 0 (if she behaves) or \(BI_i\) (if she does not).
The associated probabilities of success are \(p_H\) and \(p_L\), respectively.

If firm \(j\) is in distress, it sells its assets, which now have no internal use. We assume that
potential buyers outside the industry do not have the knowledge to operate these assets. Only firm
\(i\) (if it itself is not in distress) can buy it. There has been no initial contract that would
specify the transfer price in the case of distress. Rather, this transfer price is determined
through bargaining after distress occurs. We will later determine the per-unit transfer price \(P\).
The entrepreneur in firm \(i\) then manages \((I_1 + I_2)\) units of assets, and obtains private
benefit 0 (if she behaves) or \(B(I_1 + I_2)\) (if she does not). Similarly, the income is either 0
or \(R(I_1 + I_2)\). Probabilities of success are \(p_H\) if the entrepreneur behaves and \(p_L\) if
she does not. Firm \(i\) has just grown bigger through the acquisition.

As usual, we assume that in the absence of adverse shock (\(x = 1\)), projects are viable only if
the borrower behaves,

\[\rho_1 \equiv p_H R > 1 > p_L R + B, \qquad (14.1)\]

and we make a further assumption guaranteeing that loans are finite (even for \(x\) close to 1):

\[p_H R < 1 + \frac{p_H B}{\Delta p} \quad\text{or}\quad \rho_0 \equiv p_H\left(R - \frac{B}{\Delta p}\right) < 1. \qquad (14.2)\]

\emph{Loan agreements.} Lender \(i\) and entrepreneur \(i\) secretly sign at the start a loan
agreement specifying the amount of the loan \(I_i - A_i\) and the stake \(R_{bi}\) of entrepreneur
\(i\) in the case of success (in the absence of purchase of firm \(j\)'s assets). {[}\ldots{]}

\emph{Summary of timing} (Figure 14.1, transcribed): for each firm \(i = 1, 2\): Loan agreement
\(i\) → \(I_i\) → learning period → no distress (prob. \(x\)) or distress (prob. \(1 - x\)), with
transfer of assets from a distressed firm to the other firm → MH\(_i\) (moral hazard in firm \(i\))
→ Outcome \(i\).

\emph{Correlation of shocks.} The shocks affecting the demands for the two products may be
correlated. We allow for an arbitrary level of correlation. The conditional probabilities (given
firm \(i\)'s state) that firm \(j\) is productive or in distress are stated in Table 14.1. For
consistency, the parameters must be such that the probability that firm \(j\) is productive is
\(x\):

\[x\mu + (1 - x)(1 - \nu) = x \iff x(1 - \mu) = (1 - x)(1 - \nu). \qquad (14.3)\]

\textbf{Table 14.1}

{\def\LTcaptype{none} % do not increment counter
\begin{longtable}[]{@{}lll@{}}
\toprule\noalign{}
when firm \(i\) is & firm \(j\) productive & firm \(j\) in distress \\
\midrule\noalign{}
\endhead
\bottomrule\noalign{}
\endlastfoot
productive (prob. \(x\)) & \(\mu\) & \(1 - \mu\) \\
in distress (prob. \(1 - x\)) & \(1 - \nu\) & \(\nu\) \\
\end{longtable}
}

{[}\ldots{]} \emph{Transfer price.} If both firms are in distress (which has probability
\((1 - x)\nu\)), the four participants (entrepreneurs, lenders) receive no \emph{ex post} revenue.
If neither is in distress (which has probability \(x\mu\)), no sale occurs and the model is the
standard variable-investment one.

So let us consider the more interesting case in which \emph{firm 1, say, is in distress and firm 2
is not}. We then assume that \emph{lender 1 makes a take-it-or-leave-it offer to lender 2} (see the
third remark below for more general bargaining powers). Let \(P\) denote the per-unit price demanded
by lender 1.

Note that lender 2 must adjust entrepreneur 2's incentive scheme to account for the increased
investment and therefore for the increased private benefit from not behaving (now equal to
\(B(I_1 + I_2)\) instead of \(BI_2\)). Assume that entrepreneur 2's incentive compatibility
constraint is binding in the absence of a purchase (\((\Delta p)R_{b2} = BI_2\)), which actually
turns out to be optimal. Then, lender 2 must raise entrepreneur 2's income in the case of success by
\(\delta R_{b2}\) such that

\[(\Delta p)(\delta R_{b2}) = BI_1.\]

So, entrepreneur 2's rent increases by \([p_H B/\Delta p]I_1\) and the transfer price is

\[PI_1 = p_H\left[R - \frac{B}{\Delta p}\right]I_1 = \rho_0 I_1.\]

The transfer price is simply the pledgeable income:

\[P = \rho_0.\]

{[}\ldots{]} From (14.2), we see that \(P < 1\). {[}\ldots{]} Because firm 1 has the bargaining
power, lender 2 actually gains nothing from firm 1's distress, while entrepreneur 2 pockets the
extra agency rent \(p_H BI_1/\Delta p = (\rho_1 - \rho_0)I_1\).

\emph{Entrepreneurs' expected utility.} Suppose entrepreneur \(i\) maximizes her net utility. As
usual, the lenders' zero-profit condition implies that the borrower receives the full surplus
associated with investment:

\[U_{bi} = [xp_H RI_i - I_i] + [(1 - x)(1 - \nu)PI_i] + \left[x(1 - \mu)\frac{p_H B}{\Delta p} I_j\right]. \qquad (14.4)\]

The first term in brackets in the expression of \(U_{bi}\) corresponds to the case in which distress
sales are impossible. The second term comes from lender \(i\)'s revenue from the sales of assets if
only firm \(i\) is in distress. The third term represents the expected windfall gain from firm
\(j\)'s distress.

We can rewrite (14.4) as

\[U_{bi} = [x\rho_i + (1 - x)(1 - \nu)\rho_0 - 1]I_i + [x(1 - \mu)(\rho_1 - \rho_0)]I_j. \qquad (14.5)\]

\emph{(The book prints \(x\rho_i\); it is \(x\rho_1\), as the definition of \(\alpha\) just below
confirms.)}

So let

\[\alpha \equiv x\rho_1 + (1 - x)(1 - \nu)\rho_0 - 1\]

and

\[\kappa \equiv x(1 - \mu)(\rho_1 - \rho_0) > 0.\]

Using (14.4), we can rewrite borrower \(i\)'s net utility \(U_{bi}\) as

\[U_{bi} = \alpha I_i + \kappa I_j. \qquad (14.6)\]

\emph{Borrowing capacity.} Because lender \(i\) expects no surplus from the purchase of firm \(j\)'s
assets when the latter is in distress, lender \(i\)'s expected profit is

\[xp_H(RI_i - R_{bi}) + (1 - x)(1 - \nu)PI_i - (I_i - A_i) = 0,\]

where \(R_{bi}\) is borrower \(i\)'s income when firm \(i\) does not purchase firm \(j\)'s assets
and is successful. Incentive compatibility requires that

\[(\Delta p)R_{bi} \geqslant BI_i,\]

and, as usual, this inequality is satisfied with equality in order to maximize pledgeable income and
therefore debt capacity. Using these two equations, firm \(i\)'s maximal investment is

\[I_i = kA_i,\]

where

\[k = \frac{1}{1 - \rho_0[x + (1 - x)(1 - \nu)]}. \qquad (14.7)\]

{[}\ldots{]} We also see that the multiplier decreases with \(\nu\). That is, \emph{a firm's
borrowing capacity decreases with the degree of correlation between firms.} A higher correlation
means that the assets are less redeployable.

As we noted, the correlation of shocks reduces the desirability of investment
(\(\partial\alpha/\partial\nu < 0\)). Let \(\nu^*\) denote the level of \(\nu\) (if it exists) such
that \(\alpha = 0\). We consider two cases:

\begin{enumerate}
\def\labelenumi{(\alph{enumi})}
\item
  \emph{Low correlation} (\(\nu < \nu^*\)). Then \(\alpha > 0\) and the firms invest up to their
  borrowing capacity: \(I_i = kA_i\). While each firm's investment is independent of the existence
  of the other firm, each firm derives a positive externality from this existence in the form of
  asset redeployability.
\item
  \emph{High correlation} (\(\nu > \nu^*\)). Then \(\alpha < 0\), and no firm invests, even though
  coordinated investment could be profitable (if \(\alpha + \kappa > 0\)). {[}\ldots{]}
\end{enumerate}

\emph{Remark (alternative distributions of bargaining power).} The analysis above, and the rest of
Section 14.2 unless otherwise stated, assumes that the acquired firm's investors have full
bargaining power and thereby can charge a price equal to the acquiring investors' willingness to pay
(\(P = \rho_0\)). More generally, depending on the two parties' relative bargaining powers, the
transaction price can fall anywhere in the range between 0 (the acquired firm's opportunity cost)
and \(\rho_0\) (the acquiring firm's investors' willingness to pay).
\stmtend
%% END-STATEMENT

\asked Redo \S14.2.2 when the distressed firm's assets sell at any per-unit price $P\in[0,\rho_0]$
rather than $P=\rho_0$: the entrepreneurs' utilities (the analogues of $\alpha$ and $\kappa$), the
borrowing capacity, the interaction between the two firms' investments and the symmetric equilibrium.

\approach Tirole's continuous-investment moral-hazard model (pledgeable income $\rho_0$, NPV $\rho_1$
per unit) applied to an industry: assets of a distressed firm are worth something only to the other
firm, whose investors can pay at most the pledgeable income $\rho_0$ (the \emph{redeployability} and
\emph{fire-sale} logic). A lower $P$ moves surplus from the seller's to the buyer's investors; competitive
lending passes each lender's surplus to its entrepreneur.

\notation
\begin{itemize}
  \item $I_i$, $A_i$: firm $i$'s investment and net worth; $x$: probability a firm is productive; $\mu$, $\nu$: Table~14.1 conditional probabilities, $x(1-\mu)=(1-x)(1-\nu)$ (14.3).
  \item $\rho_1=p_HR$: per-unit NPV; $\rho_0=p_H(R-B/\Delta p)<1$: per-unit pledgeable income; $\rho_1-\rho_0=p_HB/\Delta p$: entrepreneur's per-unit rent.
  \item $P\in[0,\rho_0]$: per-unit resale price; $\hat\alpha$, $\hat\kappa$: coefficients of $U_{bi}=\hat\alpha I_i+\hat\kappa I_j$.
\end{itemize}

\step{1}{Who gets what in a transaction} When firm $j$ is distressed and $i$ productive (probability
$x(1-\mu)$), lender $i$ pays $PI_j$ and raises entrepreneur $i$'s stake so that $(\Delta p)\delta R_{bi}=BI_j$;
the acquired assets yield expected income $\rho_1I_j$, of which $(\rho_1-\rho_0)I_j$ is the entrepreneur's
rent and $\rho_0I_j$ goes to lender $i$. Per unit:
\begin{align}
  \text{seller's investors: } P, \qquad \text{buyer's investors: } \rho_0-P, \qquad \text{buyer's entrepreneur: } \rho_1-\rho_0. \label{r9e14-split}
\end{align}

\step{2}{Entrepreneur's utility} Lenders break even, so entrepreneur $i$ gets the NPV of everything her
firm does: own operations, the sale of her assets when she alone is distressed (probability $(1-x)(1-\nu)$),
and the purchase of $j$'s assets:
\begin{align}
  U_{bi} &= [x\rho_1-1]I_i+(1-x)(1-\nu)PI_i+x(1-\mu)[(\rho_1-\rho_0)+(\rho_0-P)]I_j \nonumber\\
         &= [x\rho_1+(1-x)(1-\nu)P-1]I_i+x(1-\mu)(\rho_1-P)I_j \equiv \hat\alpha I_i+\hat\kappa I_j. \label{r9e14-U}
\end{align}
Comparing with $\alpha=x\rho_1+(1-x)(1-\nu)\rho_0-1$ and $\kappa=x(1-\mu)(\rho_1-\rho_0)$,
\begin{align}
  \hat\alpha &= \alpha-(1-x)(1-\nu)(\rho_0-P), &
  \hat\kappa &= \kappa+x(1-\mu)(\rho_0-P), \label{r9e14-ak}
\end{align}
and by (14.3) the two corrections are equal, so $\hat\alpha+\hat\kappa=\alpha+\kappa$: bargaining power
only redistributes surplus between the firms for given investments.

\step{3}{Borrowing capacity} Lender $i$'s breakeven now includes its surplus $\rho_0-P$ on acquisitions:
\begin{align}
  x\rho_0I_i+(1-x)(1-\nu)PI_i+x(1-\mu)(\rho_0-P)I_j &= I_i-A_i. \label{r9e14-be}
\end{align}
Collecting the $I_i$ terms on the right,
\begin{align}
  A_i+x(1-\mu)(\rho_0-P)I_j &= I_i[1-x\rho_0-(1-x)(1-\nu)P] \nonumber\\
  I_i &= \frac{A_i+x(1-\mu)(\rho_0-P)I_j}{1-x\rho_0-(1-x)(1-\nu)P} \label{r9e14-I}\\
      &= \frac{A_i+x(1-\mu)(\rho_0-P)I_j}{1+(1-x)(1-\nu)(\rho_0-P)-\rho_0[x+(1-x)(1-\nu)]}. \nonumber
\end{align}
For $P<\rho_0$, $\partial I_i/\partial I_j=x(1-\mu)(\rho_0-P)/[1-x\rho_0-(1-x)(1-\nu)P]>0$: investments are
\emph{strategic complements} (a cheap future acquisition raises what lender $i$ will lend today), while
firm $i$'s own collateral is worth less ($P$ instead of $\rho_0$). For $P=\rho_0$ this is (14.7).

\step{4}{Symmetric equilibrium} With $A_1=A_2=A$ and $I_1=I_2=I$, \eqref{r9e14-I} gives
\begin{align}
  I[1-x\rho_0-(1-x)(1-\nu)P-x(1-\mu)(\rho_0-P)] &= A. \nonumber
\end{align}
Using $x(1-\mu)=(1-x)(1-\nu)$, the $P$ terms cancel: $-(1-x)(1-\nu)P-(1-x)(1-\nu)\rho_0+(1-x)(1-\nu)P$, so
\begin{align}
  I &= \frac{A}{1-\rho_0[x+(1-x)(1-\nu)]} = kA, \label{r9e14-sym}
\end{align}
independent of $P$: the lower collateral value and the cheaper acquisition offset each other.

\step{5}{Investment decision} Each entrepreneur invests if $\hat\alpha>0$ (her own investment's
marginal value; $\hat\kappa I_j$ is a windfall she does not control). Because $\hat\alpha\le\alpha$, the
threshold correlation $\hat\nu^*$ where $\hat\alpha=0$ is below $\nu^*$: lower seller bargaining power
shrinks the region where firms invest, although total surplus per symmetric firm $(\hat\alpha+\hat\kappa)I$
is unchanged when they do.

\step{6}{Illustration} $x=0.8$, $\nu=0.3$ ($\mu=1-0.14/0.8=0.825$), $\rho_1=1.5$, $\rho_0=0.6$, $P=0.3$, $A=1$:
$\alpha=1.2+0.084-1=0.284$, $\kappa=0.8(0.175)(0.9)=0.126$, $\hat\alpha=0.284-0.14(0.3)=0.242$,
$\hat\kappa=0.126+0.042=0.168$, and $\hat\alpha+\hat\kappa=0.41=\alpha+\kappa$. The denominator of \eqref{r9e14-I}
is $1-0.48-0.042=0.478$, the slope $0.042/0.478=0.0879$, and $I=1/(1-0.6\times0.94)=1/0.436=2.294$ for any $P$,
NPV $0.41\times2.294=0.940$. With $x=0.6$ instead, $\nu^*=1-0.1/(0.4\times0.6)=0.583$ but $\hat\nu^*=1-0.1/(0.4\times0.3)=0.167$.
\answer{(14.2)} $U_{bi}=\hat\alpha I_i+\hat\kappa I_j$ with $\hat\alpha=\alpha-(1-x)(1-\nu)(\rho_0-P)$,
$\hat\kappa=\kappa+x(1-\mu)(\rho_0-P)$, $\hat\alpha+\hat\kappa=\alpha+\kappa$;
$I_i=[A_i+x(1-\mu)(\rho_0-P)I_j]/[1-x\rho_0-(1-x)(1-\nu)P]$, increasing in $I_j$ when $P<\rho_0$; in symmetric
equilibrium $I=A/\{1-\rho_0[x+(1-x)(1-\nu)]\}$, independent of $P$.

\intuition The resale price is a transfer between the two industry members. It changes each firm's
incentive to invest (selling cheap hurts, buying cheap helps) and makes each firm's debt capacity
depend on the other's size; but in a symmetric industry every firm is equally likely to be the seller
and the buyer, so the transfers wash out of debt capacity and total value.
\pitfalls
\begin{itemize}
  \item Leaving lender $i$'s acquisition surplus $(\rho_0-P)I_j$ out of the breakeven condition; it is what makes capacity depend on $I_j$.
  \item Using $\alpha+\kappa$ to decide investment; an individual firm invests on $\hat\alpha$, the windfall term does not depend on its own choice.
\end{itemize}
\bookans Matches the book's answer (p.~618): $\hat\alpha$, $\hat\kappa$, the capacity formula and the
$P$-independence of symmetric investment.
\refbib{Tirole \S14.2.2, pp.~500--503.}

% =============================================================================
\exsec{15. Liquidity management and acquisitions --- Tirole 14.3, p.~515}
%% BEGIN-STATEMENT 15
\stmtbegin
\textbf{Source:} Tirole, Ch. 14, Exercise 14.3, p.~515 (PDF 528) · Answer in book: yes, pp.~618--619

\textbf{Exercise 14.3 (liquidity management and acquisitions).} Consider the model of Section 14.2.5
when the retooling cost is random. Suppose that this retooling cost is drawn from cumulative
distribution function \(F(\rho)\) on \([0, \infty)\), with density \(f(\rho)\) and monotonic hazard
rate (\(f(\rho)/F(\rho)\) is decreasing). The level of the retooling cost is privately observed by
the potential acquirer (the safe firm). The timing is as described in Figure 14.2.

Assume that the safe firm's entrepreneur and investors ex ante secretly agree on an investment level
\(I\) and a credit line \(L\). This credit line can be used if needed for the acquisition by the
entrepreneur and completed by the liquidity, \(\rho_0 I\), that can be raised through a seasoned
offering that dilutes the initial investors. (Fixing a credit line \(L\) of this sort is indeed an
optimal policy.)

One will assume that the seller always has the bargaining power (\(z = 1\) in the notation of
Section 14.2.5) and therefore sets price \(P\). Lastly, let \(\rho^*\) denote the equilibrium
threshold for the retooling cost (that is, assets in equilibrium are acquired and retooled if and
only if \(\rho \leqslant \rho^*\)).

\begin{enumerate}
\def\labelenumi{(\roman{enumi})}
\tightlist
\item
  Write the entrepreneur's optimal liquidity management (to this end, follow the steps described in
  Chapter 5). Show that given (anticipated) equilibrium price \(P\), the threshold \(\rho^*\)
  satisfies the ``indifference between make and buy'' equation:
\end{enumerate}

\[P + \rho^* = 1.\]

\begin{enumerate}
\def\labelenumi{(\roman{enumi})}
\setcounter{enumi}{1}
\item
  Write the objective function of the risky firm when in distress. Compute the equilibrium price
  \(P\). Note that \(P < 1\). What happens to \(P\) if for some reason the anticipated level \(L\)
  increases?
\item
  Suppose that the cumulative distribution function \(F(\rho)\) converges to a spike at
  \(\bar\rho\).\textsuperscript{35} Show that \(P + \bar\rho = 1\), and that \(F(\rho^*)\) converges
  to 1.
\end{enumerate}

\textsuperscript{35} While still satisfying the monotone hazard rate property.

\paragraph{Context needed to solve it}

\emph{From Tirole §14.2.2, p.~500 (PDF 513) --- per-unit NPV and pledgeable income:}

\[\rho_1 \equiv p_H R > 1 > p_L R + B, \qquad (14.1)\]

\[p_H R < 1 + \frac{p_H B}{\Delta p} \quad\text{or}\quad \rho_0 \equiv p_H\left(R - \frac{B}{\Delta p}\right) < 1. \qquad (14.2)\]

\emph{From Tirole §14.2.5.3 ``Bilateral Monopoly: Is Liquidity Hoarding Held Up?'', pp.~506--508
(PDF 519--521):}

Let us first describe the model in the case of a single risky firm and a single safe firm. The two
firms choose investment sizes \(J\) and \(I\), respectively. To simplify notation without altering
the basic insights, the two firms are identical, except for the probability of distress, which is 0
for the safe firm and \(1 - x\) for the risky one. Both are run by risk-neutral entrepreneurs with
initial cash on hand \(A\) each and protected by limited liability. The timing is summarized in
Figure 14.2.

\emph{Figure 14.2 (transcribed).} Risky firm: invests \(J\) → learns whether productive (prob.
\(x\)) → moral hazard: behaves (\(p_H\), no private benefit) or misbehaves (\(p_L\), private benefit
\(BJ\)) → outcome: success (\(RJ\)) or failure (0). With prob. \(1 - x\): assets resale to the safe
firm. Safe firm: invests \(I\), hoards liquidity \(L\) → (if acquisition) pays retooling cost
\(\rho J\) (on top of acquisition cost \(PJ\)) → moral hazard: behaves (\(p_H\), no private benefit)
or misbehaves (\(p_L\), private benefit \(BI\) or \(B(I + J)\)) → outcome: success (\(RI\) or
\(R(I + J)\)) or failure (0).

We further assume that when buying the \(J\) units of asset from the risky firm when the latter is
in distress, the safe firm must pay a known retooling cost \(\rho J\) to adapt these assets to its
own production process. Letting \(P\) denote the unit acquisition price, the total cost of the
acquisition is thus \((P + \rho)J\). To limit the number of cases to be considered, we assume that
\(\rho \leqslant \rho_0\).

{[}\ldots{]} As we want to allow for a wide range of bargaining powers, let us assume that in the
case of distress of the risky firm:

\begin{itemize}
\tightlist
\item
  with probability \(z\), the risky firm makes a take-it-or-leave-it resale offer to the safe firm;
\item
  with probability \(1 - z\), the safe firm makes a take-it-or-leave-it purchase offer.
\end{itemize}

Thus \(z\) is a measure of the selling firm's bargaining power.

\emph{Bargaining and the choice of financial muscle.} The case in which the safe firm makes the
offer is a no-brainer: it offers 0 (or just above), the opportunity cost of the risky firm's assets
when in distress. By contrast, the situation in which the risky firm makes the offer requires more
thinking. Let us assume that the amount of liquidity \(L\) hoarded by the safe firm (say, the credit
line that the firm secures from its bank\textsuperscript{24}) is \emph{not} observed by the risky
firm. We look for conditions under which the following is an equilibrium:

\begin{itemize}
\tightlist
\item
  the risky firm demands per-unit price \(\rho_0\) for its assets;
\item
  the safe firm, anticipating this and knowing that it will be able to raise \(\rho_0 J\) on the
  asset market, hoards liquidity \(L = \rho J\).
\end{itemize}

Making offer \(\rho_0\) is clearly optimal for the seller since \(\rho_0 J\) is the pledgeable
income on the acquired assets and thus the upper bound on what new investors are willing to
contribute in a seasoned offering.

Let us next investigate whether it is indeed in the safe firm's interest to hoard liquidity. Let
\(y = 1\) if it hoards the necessary amount (\(\rho J\)) and \(y = 0\) if it does not. Noting that
liquidity is not needed to acquire the assets when the acquirer makes the offer (since
\(\rho \leqslant \rho_0\)), the safe firm's NPV is then

\[U_b^s = (\rho_1 - 1)I + (1 - x)[z[\rho_1 - (\rho_0 + \rho)]y + (1 - z)(\rho_1 - \rho)]J.\]

The second term on the right-hand side is the expected gain from acquisitions, and uses the fact
that the total acquisition cost is \((\rho_0 + \rho)J\) when the seller makes the offer and only
\(\rho J\) when the acquirer makes the offer.

The safe firm's investors must break even, and so

\[\rho_0 I + (1 - x)[z[\rho_0 - (\rho_0 + \rho)]y + (1 - z)(\rho_0 - \rho)]J = I - A.\]

The investors' breakeven condition yields the investment level \(I\):

\[I = \frac{A + [(1 - z)(\rho_0 - \rho) - z\rho y](1 - x)J}{1 - \rho_0}.\]

Note that hoarding liquidity (\(y = 1\)) reduces the investment scale. Substituting into the NPV
equation,

\[U_b^s = (\rho_1 - 1)\frac{A + [(1 - z)(\rho_0 - \rho) - z\rho y](1 - x)J}{1 - \rho_0} + (1 - x)[z[\rho_1 - (\rho_0 + \rho)]y + (1 - z)(\rho_1 - \rho)]J.\]

We conclude that it is an equilibrium for the safe firm to hoard liquidity to purchase assets if and
only if \(U_b^s\) is increasing in \(y\), or

\[\rho_1 - (\rho_0 + \rho) \geqslant \frac{(\rho_1 - 1)\rho}{1 - \rho_0}\]

or

\[1 \geqslant \rho_0 + \rho.\]

We thus obtain the simple result that the potential acquirer builds the necessary financial muscle
if and only if the total per unit cost of acquisition when the acquired firm has bargaining power is
lower than the safe firm's own cost of investment. {[}\ldots{]}

In this simple model the equilibrium is in general not unique: there are lots of other
self-fulfilling equilibria in which the firm hoards \(L^* \neq \rho J\) and the seller demands
\(P^*\) such that \(P^* + \rho = L^* + \rho_0\). The seller does not want to demand more than
\(P > P^*\) because being too greedy prevents the potential acquirer from buying the assets.
Conversely, the potential acquirer is willing to hoard \(L = L^*\), as long as
\(P^* + \rho \leqslant 1\) (note that \(\rho_0\) is in the interior of the range of equilibrium
prices).

\textsuperscript{24} We adopt the convention that the safe firm's entrepreneur can raise \(\rho_0\)
in a seasoned security offering, to which \(L\) is added to form total available cash.
Alternatively, and as we discussed in Chapter 5, one could ban seasoned offerings and provide the
firm with a bigger credit line.

\emph{From Tirole §5.3.1 ``The Two-Shock Case'', pp.~207--209 (PDF 220--222) --- ``the steps
described in Chapter 5'':}

{[}\ldots{]} (ii) \emph{Pursue the project even in the case of distress.} The decision to withstand
the liquidity shock at date 1 has a cost and benefit. The cost is that the average cost of bringing
1 unit of investment intact to date 2 is \((1 + \lambda\rho)\) (the date-0 cost, 1, plus the
expected date-1 reinvestment cost, \(\lambda\rho\)). The benefit is that the project is never
abandoned. The borrowing capacity is given by

\[(1 + \lambda\rho)I - A = \rho_0 I \quad\text{or}\quad I = \frac{A}{(1 + \lambda\rho) - \rho_0}.\]

Similarly, the entrepreneur's utility (the NPV) is

\[U_b^1 = [\rho_1 - (1 + \lambda\rho)]I \quad\text{or}\quad U_b^1 = \frac{\rho_1 - (1 + \lambda\rho)}{(1 + \lambda\rho) - \rho_0}A. \quad […]\]

{[}\ldots{]} If the optimal policy is to pursue the project even in the case of distress, the
entrepreneur must be able to avail herself of the amount \(\rho I\) if a shock occurs. If
\(\rho > \rho_0\) {[}\ldots{]}, then liquidity necessarily must be planned in advance. Waiting
exposes the firm to credit rationing at date 1. {[}\ldots{]} For example, the firm may contract a
credit line to the level of \(\rho I\) with a bank; alternatively, it can contract for a credit line
corresponding only to the shortfall \((\rho - \rho_0)I\) and also acquire the right to dilute
initial investors (so as to obtain \(\rho_0 I\)).
\stmtend
%% END-STATEMENT

\asked (i) The safe firm's optimal credit line and acquisition cutoff when the retooling cost $\rho$ is
random and privately observed, given the price $P$: show $P+\rho^*=1$. (ii) The distressed seller's
pricing problem, the equilibrium $P<1$, and how $P$ responds to a larger credit line. (iii) The limit
when $F$ collapses to a point $\bar\rho$.

\approach Chapter~5's liquidity--scale trade-off: a unit of liquidity reserved for a future need costs
investment scale today (multiplier $1/(1-\rho_0)$), and a credit line caps which shocks are met. Here the
``shock'' is an acquisition opportunity of cost $(P+\rho)J$. The seller is a monopolist facing an
unknown buyer's maximum willingness to pay, so it trades off price against the probability of a sale
(monotone hazard rate gives a unique price).

\notation
\begin{itemize}
  \item $I$, $J$: safe and risky firms' investments; $A$: safe firm's net worth; $1-x$: probability the risky firm is distressed.
  \item $\rho\sim F$ on $[0,\infty)$, density $f$, $f/F$ decreasing; $G(\rho)\equiv F(\rho)/f(\rho)$, so $G'>0$.
  \item $L$: credit line per unit of $J$; $\rho_0J$: cash from a seasoned offering; the acquisition is feasible iff $(P+\rho)J\le(\rho_0+L)J$, so $\rho^*=\rho_0+L-P$.
\end{itemize}

\step{1}{(i) NPV and breakeven} The acquisition, when it happens, yields NPV $\rho_1-(P+\rho)$ per unit
and pledgeable income $\rho_0-(P+\rho)$ per unit to investors (who fund it through the line and the
offering). Following Chapter~5,
\begin{align}
  U_b^s &= (\rho_1-1)I+(1-x)J\int_0^{\rho^*}[\rho_1-(P+\rho)]\,dF(\rho), \label{r9e15-U}\\
  \rho_0I+(1-x)J\int_0^{\rho^*}[\rho_0-(P+\rho)]\,dF(\rho) &= I-A. \label{r9e15-be}
\end{align}
Solving \eqref{r9e15-be} for $I$ (move $\rho_0I$ to the right and divide by $1-\rho_0>0$):
\begin{align}
  I &= \frac{A-(1-x)J\int_0^{\rho^*}[(P+\rho)-\rho_0]\,dF(\rho)}{1-\rho_0}. \label{r9e15-I}
\end{align}
Substituting \eqref{r9e15-I} in \eqref{r9e15-U}:
\begin{align}
  U_b^s &= (\rho_1-1)\frac{A-(1-x)J\int_0^{\rho^*}[(P+\rho)-\rho_0]\,dF}{1-\rho_0}+(1-x)J\int_0^{\rho^*}[\rho_1-(P+\rho)]\,dF. \label{r9e15-U2}
\end{align}

\step{2}{(i) Optimal cutoff} Differentiate \eqref{r9e15-U2} in $\rho^*$ (Leibniz rule: each integral contributes
its integrand at $\rho^*$ times $f(\rho^*)$), and write $s=P+\rho^*$:
\begin{align}
  \frac{\partial U_b^s}{\partial\rho^*} &= (1-x)Jf(\rho^*)\Big[-\frac{(\rho_1-1)(s-\rho_0)}{1-\rho_0}+\rho_1-s\Big]. \nonumber
\end{align}
Setting the bracket to zero and multiplying by $1-\rho_0$:
\begin{align}
  (\rho_1-s)(1-\rho_0) &= (\rho_1-1)(s-\rho_0) \nonumber\\
  \rho_1-\rho_1\rho_0-s+s\rho_0 &= \rho_1s-\rho_1\rho_0-s+\rho_0 \nonumber\\
  \rho_1+s\rho_0 &= \rho_1s+\rho_0 \nonumber\\
  \rho_1-\rho_0 &= s(\rho_1-\rho_0) \nonumber\\
  s &= 1. \nonumber
\end{align}
The bracket is positive for $s<1$ and negative for $s>1$ (it is linear in $s$ with slope
$-(\rho_1-1)/(1-\rho_0)-1<0$), so this is a maximum:
\begin{align}
  P+\rho^* &= 1, \qquad\text{credit line } L^*=\rho^*+P-\rho_0=1-\rho_0 \text{ per unit of } J. \label{r9e15-cut}
\end{align}
\answer{(i)} The safe firm secures a credit line $L^*J$ with $\rho_0+L^*=P+\rho^*=1$: it acquires iff the
total unit cost of buying, $P+\rho$, is at most 1, the unit cost of growing internally (``make or buy'').
Numerically, with $F$ exponential ($\lambda=2$), $\rho_0=0.6$, $\rho_1=1.5$, $x=0.8$, $J=A=1$ and $P=0.4$, a grid
search over $\rho^*$ maximizes $U_b^s$ at $\rho^*=0.6=1-P$.

\step{3}{(ii) Seller's problem} The seller does not observe $\rho$ and anticipates the line $L$; a sale
occurs iff $\rho\le\rho_0+L-P$. Its assets have no other use, so it maximizes expected revenue per unit:
\begin{align}
  \max_P\; P\,F(\rho_0+L-P). \label{r9e15-seller}
\end{align}
The first-order condition (with $\rho^*=\rho_0+L-P$) is
\begin{align}
  F(\rho^*)-Pf(\rho^*) &= 0 \iff P = \frac{F(\rho^*)}{f(\rho^*)}=G(\rho^*). \label{r9e15-foc}
\end{align}
In equilibrium $L=L^*=1-\rho_0$, so $\rho^*=1-P$ and $P$ solves $\phi(P)\equiv -P+G(1-P)=0$. At $P=0$,
$\phi=G(1)>0$; at $P=1$, $\phi=-1+G(0)=-1<0$ (as $F(0)=0$); and $\phi'(P)=-1-G'(1-P)<0$ because $G'>0$ by the
monotone hazard rate. Hence a unique $P\in(0,1)$.

\step{4}{(ii) Larger credit line} Differentiate $-P+G(\rho_0+L-P)=0$ in $L$:
\begin{align}
  -\frac{dP}{dL}+G'\Big(1-\frac{dP}{dL}\Big) &= 0 \;\Rightarrow\; \frac{dP}{dL}=\frac{G'}{1+G'}\in(0,1). \label{r9e15-dPdL}
\end{align}
The cutoff rises, $d\rho^*/dL=1-dP/dL=1/(1+G')>0$: the seller captures part of a larger line but the
probability of a sale still increases. Example $F(\rho)=1-e^{-2\rho}$: $G(\rho)=(e^{2\rho}-1)/2$,
$G'=e^{2\rho}$; the equilibrium is $P=0.6040$, $\rho^*=0.3960$, $F(\rho^*)=0.547$, and
$dP/dL=e^{0.792}/(1+e^{0.792})=0.688$ (confirmed by finite differences).
\answer{(ii)} The risky firm maximizes $PF(\rho_0+L-P)$; the equilibrium price solves $P=F(1-P)/f(1-P)$,
unique and in $(0,1)$; a higher anticipated $L$ raises $P$ with pass-through $G'/(1+G')<1$, so the
probability of a deal rises too.

\step{5}{(iii) Spike at $\bar\rho$} Informally (as in the book): a price whose cutoff $1-P$ stays below
$\bar\rho$ sells with probability near 0, so the seller does better by setting a cutoff just above
$\bar\rho$; a cutoff that stays above $\bar\rho$ gives $Pf\approx0$ and $F\approx1$, violating the
first-order condition $F=Pf$ unless the cutoff converges to $\bar\rho$. So $1-P\to\bar\rho$ and $F(\rho^*)\to1$.
Explicit sequence: $F$ uniform on $[\bar\rho-\varepsilon,\bar\rho+\varepsilon]$ ($f/F=1/(\rho-\bar\rho+\varepsilon)$
decreasing). Revenue $P\cdot\min\{(1-P-\bar\rho+\varepsilon)/(2\varepsilon),1\}$ increases in the cutoff up to
$\rho^*=\bar\rho+\varepsilon$ when $1-\bar\rho>3\varepsilon$, so $P=1-\bar\rho-\varepsilon$ and $F(\rho^*)=1$; as
$\varepsilon\to0$, $P\to1-\bar\rho$. Example $\bar\rho=0.5$, $\varepsilon=0.05$: $P=0.45$, sale with probability 1.
\answer{(iii)} $P+\bar\rho=1$ in the limit and $F(\rho^*)\to1$: the price equals the buyer's make-or-buy
indifference price, the competitive outcome, and a deal is (almost) sure.

\intuition The buyer hoards exactly the liquidity that makes buying as cheap as building; the seller,
knowing this, prices like a monopolist against the buyer's uncertain retooling cost. With little
uncertainty the monopolist extracts everything up to the make-or-buy price and the deal always
happens; with more uncertainty it prices some deals away, so trade is inefficiently low.
\pitfalls
\begin{itemize}
  \item Forgetting the opportunity cost of liquidity: hoarding reduces $I$ through \eqref{r9e15-I}; ignoring it gives the wrong cutoff $\rho^*=\rho_1-P$.
  \item Treating $L$ as observed by the seller out of equilibrium; the seller prices on the \emph{anticipated} line.
\end{itemize}
\bookans Matches the book's answer (pp.~618--619): $\rho^*=1-P$, $\rho_0+L^*=1$; $P$ maximizes $F(\rho_0+L^*-P)P$,
unique in $(0,1)$; $0<dP/dL<1$; informal limit argument for (iii).
\refbib{Tirole \S14.2.5.3, pp.~506--508; \S5.3.1, pp.~207--209.}

\end{document}
